\documentclass[11pt, a4paper]{amsart}
\usepackage{amsmath, amssymb, amsthm, amsfonts}
\usepackage{mathtools}
\usepackage[margin=1in]{geometry}
\usepackage{booktabs}
\usepackage{hyperref}
\usepackage{tikz-cd}

\newtheorem{theorem}{Theorem}[section]
\newtheorem{lemma}[theorem]{Lemma}
\newtheorem{proposition}[theorem]{Proposition}
\newtheorem{corollary}[theorem]{Corollary}
\newtheorem{fact}[theorem]{Standard Fact}
\theoremstyle{definition}
\newtheorem{definition}[theorem]{Definition}
\newtheorem{remark}[theorem]{Remark}
\newtheorem{notation}[theorem]{Notation}
\newtheorem{example}[theorem]{Example}

\newcommand{\R}{\mathbb{R}}
\newcommand{\Q}{\mathbb{Q}}
\newcommand{\C}{\mathbb{C}}
\newcommand{\Z}{\mathbb{Z}}
\newcommand{\HH}{\mathbb{H}}
\newcommand{\OO}{\mathbb{O}}
\newcommand{\id}{\operatorname{id}}
\newcommand{\Der}{\operatorname{Der}}
\newcommand{\Aut}{\operatorname{Aut}}
\newcommand{\im}{\operatorname{im}}
\newcommand{\Stab}{\operatorname{Stab}}
\newcommand{\tr}{\operatorname{tr}}
\newcommand{\Cl}{\mathrm{Cl}}
\newcommand{\Spin}{\mathrm{Spin}}
\newcommand{\Sp}{\mathrm{Sp}}
\newcommand{\SO}{\mathrm{SO}}
\newcommand{\SU}{\mathrm{SU}}
\newcommand{\Orth}{\mathrm{O}}
\renewcommand{\Im}{\mathrm{Im}}

\title{Cascade--Classical Correspondence II:\\
  The Cascade Algebraic Substrate}
\author{}
\date{}

\begin{document}

\maketitle

\begin{abstract}
This paper constructs the algebraic substrate underlying the cascade
programme: the classical $E_8$ root system, the strict
sub-root-system inclusion $D_4 \subset E_8$, the icosian projection
from $E_8$ lattice data to an $H_4$ module, an explicit $24$-cell
inside the $600$-cell, the spectral and shell structure of the
$600$-cell vertex graph (its identification with the binary
icosahedral group $2I$ under Hamilton multiplication, the
Euclidean-shell-to-conjugacy-class bijection, the centrality of
every shell-adjacency class sum in $\C[2I]$ and the resulting
isotypic block structure (with the Laplacian spectrum of
Fact~\ref{prop:spectrum-P2} imported from the standard $2I$
character table to label the integer / irrational $94{+}26$
isotypic decomposition), the BFS shell structure with its
icosahedral /
dodecahedral subpopulations, the lumpable $9$-state Markov chain
on that partition with its exact rational first-passage identity
$R_D(4) = 15$, a narrow uniqueness property among the eleven
regular $3$-/$4$-polytopes, and the Hopf decomposition into
twelve great-decagon fibres), the real Clifford algebra $\Cl(1,3)$
as the canonical $16$-sector model, and the octonionic $8$-sector
with its derivation algebra $\mathfrak{g}_2 = \Der(\OO)$ and the
$\SU(3) \subset G_2$ stabilizer. The coefficientwise $\sigma$-action
on all $\Q(\varphi)$-defined objects is imported from
\cite{SigmaRationality} and shown to intertwine every $\Q$-linear
sector-level map (Theorem~\ref{thm:sigma-compat-P2}). The
non-$\Q$-linear projection $\pi_H$ does not commute with
$\sigma$; the relevant fact for $\eta_\Q$ instead is the
swap-intertwining (Proposition~\ref{prop:eta-swap}).

This is the second of thirteen planned infrastructure papers
(\cite[P2]{CascadeInfrastructurePlan}). Its purpose is to replace
informal and slogan-level references to the cascade sector sequence
$E_8 \to H_4 \to 24 \to D_4 \to 16 \to 8$ with theorem-grade
constructions supported either by classical mathematics or by
\cite{SigmaRationality}. \emph{No seven-rung uniqueness theorem is
proved here.} In particular, the paper makes no claim that the
cascade chain is forced by $E_8$ maximality, by Coxeter
classification, or by any other external principle.

The paper's main proved results (all other content is standard and
cited): (i) the strict inclusion
$\Phi(D_4) = \Phi(E_8) \cap U$ for the coordinate $4$-plane
$U = \operatorname{span}_\R(e_1, e_2, e_3, e_4)$
(Proposition~\ref{prop:D4-inclusion}); (ii) the identification of an
explicit $24$-vertex subset $V_{24} \subset V_{600}$ with a regular
$24$-cell isometric to the $D_4$ root polytope
(Proposition~\ref{prop:24-in-600}); (iii) the computer-assisted
identification of $V_{600}$ under Hamilton multiplication with the
binary icosahedral group $2I$
(Theorem~\ref{thm:icosian-closure}), together with the exact
$\Q(\sqrt 5)$-arithmetic bijection between the nine Euclidean
distance shells from the identity and the nine $2I$-conjugacy
classes (Theorem~\ref{thm:shell-class}), and the centrality of
every shell-adjacency class sum under the resulting left-regular
realisation $\C^{V_{600}} \cong \C[2I]$
(Corollary~\ref{cor:class-sums-central}); (iv) the equitable
nine-cell BFS partition with full integer quotient matrix
(Proposition~\ref{prop:dual-populations}), the strong lumpability
of the uniform walk to a $9$-state chain $\widehat P = M / 12$
(Proposition~\ref{prop:lumped-markov}), and the exact rational
first-passage identity $R_D(4) = 15$
(Theorem~\ref{thm:RD4-P2}); (v) the narrow uniqueness theorem
among the eleven regular $3$-/$4$-polytopes, conditional on the
imported classical Laplacian spectra of the ten non-$600$-cell
members (Theorem~\ref{thm:polytope-uniqueness}); the imported
Laplacian spectrum of the $600$-cell vertex graph itself
(Fact~\ref{prop:spectrum-P2}) and the resulting $94{+}26$
isotypic labelling (Corollary~\ref{cor:94-26-split}) are recorded
but rest on the imported $2I$ character table rather than on
additional P2 work; and (vi) the
$\sigma$-compat\-ibility theorem that every $\Q$-defined sector-level
map commutes with the coefficientwise $\sigma$-action after scalar
extension (Theorem~\ref{thm:sigma-compat-P2}).

All other statements are well-known classical results used through
standard references: Humphreys and Bourbaki for root systems,
Coxeter for regular polytopes, Conway--Sloane and Moody--Patera for
icosians, Lounesto for Clifford algebras, Baez for octonions and
$G_2$. We do not rely on any cascade-derivation internal note as
theorem input.
\end{abstract}

\tableofcontents

%=====================================================================
\section{Introduction and Source Audit}
\label{sec:intro}
%=====================================================================

\subsection{What this paper proves}

This paper constructs an algebraic substrate suitable for use as a
theorem-grade dependency in downstream cascade infrastructure
papers (P3--P13 of \cite{CascadeInfrastructurePlan}) and in the
cascade-correspondence Foundations paper. The substrate consists of
seven distinct algebraic objects, each with a precise model:
\begin{enumerate}
\item the $E_8$ root system and root lattice (\S\ref{sec:E8});
\item the strict sub-root-system $D_4 \subset E_8$ via the
  coordinate $4$-plane $U$ (\S\ref{sec:D4});
\item the icosian ring $I$ and the explicit projection $\pi_H$
  from $E_8$ lattice data to an $H_4$ module (\S\ref{sec:H4});
\item an explicit $24$-cell as a subset of the $600$-cell
  (\S\ref{sec:24-600});
\item the shell, representation-theoretic, and lumpability
  structure of the $600$-cell vertex graph
  (\S\ref{sec:600cell-spectral}): the identification of the $120$
  vertices with the binary icosahedral group $2I$ under Hamilton
  multiplication, and the exact $\Q(\sqrt 5)$-arithmetic bijection
  between the nine Euclidean distance shells from the identity
  and the nine $2I$-conjugacy classes; the resulting centrality
  of every shell-adjacency class sum in $\C[2I]$ (the $94{+}26$
  integer / irrational labelling of isotypic components and the
  explicit Laplacian spectrum are imported from the standard $2I$
  character table, not re-proved here); the breadth-first-search
  shell structure with its dual $(\mathrm{ico}, \mathrm{dod})$
  subpopulations with an explicit $9 \times 9$ equitable-partition
  quotient matrix; the lumped $9$-state Markov chain on that
  equitable partition, together with the absorbing first-passage
  reduction yielding the exact rational identity $R_D(4) = 15$; a narrow uniqueness property among the
  eleven regular $3$- and $4$-polytopes (conditional on imported
  Laplacian spectra for the ten non-$600$-cell members); and the
  Hopf decomposition into twelve great-decagon fibres (imported);
\item the real Clifford algebra $\Cl(1,3)$ as the algebraic model of
  the $16$-sector (\S\ref{sec:16-sector});
\item the octonionic $8$-sector with its derivation algebra
  $\Der(\OO) = \mathfrak{g}_2$ and the $\SU(3) \subset G_2$
  stabilizer of a fixed imaginary unit (\S\ref{sec:8-sector}).
\end{enumerate}

Section~\ref{sec:sigma-compat} then imports the coefficientwise
$\sigma$-action of \cite{SigmaRationality} and shows that every map
defined in \S\S\ref{sec:D4}--\ref{sec:8-sector} intertwines with
$\sigma$ in the sense of that paper.

\subsection{What this paper does not prove}

P2 is deliberately narrow. To forestall any reader from reading
this as a seven-rung uniqueness paper, we record the exclusions
explicitly:

\begin{enumerate}
\item P2 does not claim the full seven-rung chain
  $E_8 \to H_4 \to 40 \to D_4 \to 16 \to 8 \to 0$ as a theorem; no
  uniqueness, maximality, or completeness statement is made about
  the chain.
\item P2 does not claim that $H_4$ is a strict sub-root-system of
  $E_8$. $H_4$ is non-crystallographic, and we always describe it
  as the image of a projection, never as an inclusion into
  $\Phi(E_8)$.
\item P2 does not claim any vertex-wise map $V_{600} \to V_{24}$
  (in particular, there is no ``$120 \to 24$'' quotient map in this
  paper). The $24$-cell appears as an embedded subset of the
  $600$-cell, not as a quotient.
\item P2 does not use the $2I/2T$ coset combinatorics of earlier
  cascade notes as a theorem input. Any combinatorial identity
  involving icosahedral rotation groups in this paper is proved
  directly or cited to a classical source.
\item P2 does not claim that the $16$-sector is uniquely forced
  from $D_4$; $\Cl(1,3)$ is simply a natural $16$-dimensional real
  algebra, chosen for its relevance to downstream applications.
\item P2 does not claim that the $8$-sector is uniquely forced
  from triality; the octonion algebra is a natural $8$-dimensional
  normed division algebra with an $\SU(3)$-stabilizer structure,
  chosen for its relevance to later gauge-sector work in P8.
\item P2 does not use $\mathrm{ad} \colon \Im(\OO) \to
  \mathfrak{g}_2$ as a Lie-algebra inclusion. This ``ad-map''
  identification is algebraically false in general, as noted in the
  iteration-1 review of the earlier Foundations draft. We use
  derivations $D_{a,b}$ instead, which is the standard and correct
  construction \cite{BaezOctonions}.
\item P2 asserts no physical or phenomenological identification
  between these algebraic sectors and downstream gauge, metric, or
  cohomological objects. Those identifications are deferred to
  dedicated downstream papers.
\item P2 does not certify the individual numerical entries of the
  $2I$ character table by bare hand. The \emph{centrality}
  of each shell-adjacency class sum on each isotypic component of
  $\C[2I]$ (\S\ref{sec:600cell-spectral}) is proved here by
  Schur's lemma after a finite exact-arithmetic verification that
  each Euclidean distance shell is a single $2I$-conjugacy class;
  the \emph{specific numerical scalar} by which each class sum
  acts on each isotypic component is read from the standard $2I$
  character table~\cite[Ch.~3]{DuVal1964}. The
  labelling of isotypic components as ``integer-eigenvalue'' or
  ``irrational-eigenvalue'' in the $94{+}26$ decomposition
  therefore rests on that imported character data, not on
  additional representation-theoretic work in this paper.
\item P2 does not re-prove the classical Laplacian spectra of the
  non-$600$-cell members of the $11$-polytope family used in the
  narrow-uniqueness theorem of \S\ref{sec:600cell-spectral};
  those spectra are imported from \cite{BrouwerHaemers2012} and
  \cite{GodsilRoyle2001} and re-verified computationally by the
  accompanying script, but their classical derivations are not
  repeated here.
\item P2 does not prove the Hopf decomposition of the $600$-cell
  into twelve great-decagon fibres; that decomposition is
  imported from the classical $600$-cell literature and used in
  \S\ref{sec:600cell-spectral} to state the fibre cycle spectrum.
\end{enumerate}

\subsection{Source audit}

Theorem inputs to this paper split into two classes:
\begin{description}
\item[\textbf{Classical mathematics.}] The $E_8$ root system and
  lattice \cite{Humphreys, Bourbaki}; the classical theory of
  regular polytopes and the $24$- and $600$-cells
  \cite{CoxeterRegularPolytopes}; icosians and the $E_8$ lattice
  \cite{ConwaySloane, MoodyPatera}; real Clifford algebras
  \cite{Lounesto, Chevalley}; octonions and $G_2$
  \cite{BaezOctonions, SchaferNonassociative}; the Du~Val / ADE
  classification of the finite subgroups of $\mathrm{Sp}(1)$
  \cite{DuVal1964}; spectra of distance-regular graphs and of
  the Platonic / regular-polytope vertex graphs
  \cite{BrouwerHaemers2012, GodsilRoyle2001}; absorbing Markov
  chains and equitable partitions \cite{KemenySnell1976,
  GodsilRoyle2001}; and the standard $2I$ character table, as
  part of Du~Val's classification of the binary polyhedral groups
  \cite[Ch.~3]{DuVal1964}.
\item[\textbf{Internal dependency.}] Paper P1
  (\cite{SigmaRationality}), cited only for the coefficientwise
  $\sigma$-action on $\Q(\varphi)$-scalar-extended vector spaces
  and its functoriality theorem.
\end{description}

We do \emph{not} rely on any cascade-derivation internal note as
theorem input. In particular, \cite[\S F3]{CascadeFoundations}
(seven-rung uniqueness), \cite[\S\S P4--P8]{CascadePregeometry}
(bootstrap chain), and \cite[\S 12]{CascadeEmbeddings}
(computational verifications) are not used in any proof. These
sources appear in the bibliography only to record their exclusion,
and any reader who encounters a formal-framework or computational
note from those sources should not treat it as a citation of
P2.

\subsection{Organization}

Section~\ref{sec:E8} sets up the standard $E_8$ root system.
Section~\ref{sec:D4} proves the strict $D_4 \subset E_8$ inclusion.
Section~\ref{sec:H4} constructs the icosian ring and the $E_8 \to
H_4$ projection. Section~\ref{sec:24-600} isolates the $24$-cell
inside the $600$-cell by explicit coordinates.
Section~\ref{sec:600cell-spectral} collects the
spectral, shell, and representation-theoretic structure of the
$600$-cell vertex graph: the identification with the binary
icosahedral group $2I$, the Euclidean-shell-to-conjugacy-class
bijection, the Laplacian spectrum with its isotypic $94{+}26$
integer/irrational split, the BFS shell structure with its dual
$(\mathrm{ico}, \mathrm{dod})$ subpopulations, the narrow
uniqueness property among the eleven regular $3$- and
$4$-polytopes, and the Hopf decomposition into twelve great-decagon
fibres. Section~\ref{sec:16-sector} presents the $\Cl(1,3)$ model
of the $16$-sector. Section~\ref{sec:8-sector} presents the
octonionic $8$-sector. Section~\ref{sec:sigma-compat} establishes
the $\sigma$-compatibility theorem. Section~\ref{sec:handoff}
gives the citation contract for downstream papers.
Appendix~\ref{app:notation} is a notation concordance.

%=====================================================================
\section{Standard $E_8$ Data and Normalizations}
\label{sec:E8}
%=====================================================================

This section fixes conventions for the $E_8$ root system and
lattice. The content is entirely classical and is cited rather than
reproved.

\subsection{The $E_8$ root system}

\begin{definition}[$E_8$ root system]
\label{def:E8}
The \emph{$E_8$ root system} $\Phi(E_8) \subset \R^8$ consists of
the following $240$ vectors:
\begin{itemize}
\item \textbf{$D_8$-type roots ($112$):}
  $\pm e_i \pm e_j$ for $1 \leq i < j \leq 8$, with independent
  signs; there are $4 \binom{8}{2} = 112$ such vectors.
\item \textbf{Half-integer roots ($128$):}
  $\frac{1}{2}(\pm 1, \pm 1, \pm 1, \pm 1, \pm 1, \pm 1, \pm 1, \pm 1)$
  with the constraint that the number of minus signs is even;
  there are $2^7 = 128$ such vectors.
\end{itemize}
Here $\{e_1, \ldots, e_8\}$ is the standard orthonormal basis of
$\R^8$ with inner product $\langle \cdot, \cdot \rangle$.
\end{definition}

\begin{definition}[$E_8$ root lattice]
\label{def:lambda-E8}
The \emph{$E_8$ root lattice} $\Lambda(E_8) \subset \R^8$ is the
$\Z$-linear span of $\Phi(E_8)$. Equivalently,
\[
  \Lambda(E_8) = \big\{ (x_1, \ldots, x_8) \in \R^8 :
    \text{either all } x_i \in \Z \text{ or all } x_i \in \Z + \tfrac{1}{2},
    \; \sum x_i \in 2\Z \big\}.
\]
$\Lambda(E_8)$ is even, unimodular, and self-dual
\cite[\S III.4]{Serre}.
\end{definition}

\begin{fact}[Standard, \cite{Humphreys, Bourbaki}]
\label{fact:E8-irred}
$\Phi(E_8)$ is a reduced, irreducible crystallographic root system
of rank $8$ and order $240$. Its Weyl group $W(E_8)$ has order
$696\,729\,600 = 2^{14} \cdot 3^5 \cdot 5^2 \cdot 7$.
\end{fact}

\begin{remark}[Normalization]
\label{rmk:normalization}
With the conventions of Definition~\ref{def:E8}, all roots have
squared length $2$. We retain this normalization throughout. Where
other cascade sources (e.g., \cite{CascadeEmbeddings}) use
$\frac{1}{2}\|r\|^2 = 1$, the two conventions agree.
\end{remark}

%=====================================================================
\section{The Strict $D_4$ Sub-root-system of $E_8$}
\label{sec:D4}
%=====================================================================

This section proves the first load-bearing proposition of the
paper: $D_4$ embeds as a strict sub-root-system of $E_8$ via the
coordinate $4$-plane $U$.

\subsection{The coordinate $4$-plane $U$}

\begin{definition}[The plane $U$]
\label{def:U}
Let
\[
  U := \operatorname{span}_\R(e_1, e_2, e_3, e_4) \subset \R^8.
\]
$U$ is a $4$-dimensional $\R$-subspace of $\R^8$, identified with
$\R^4$ by the map $(x_1, x_2, x_3, x_4) \leftrightarrow
(x_1, x_2, x_3, x_4, 0, 0, 0, 0)$.
\end{definition}

\begin{definition}[$D_4$ root system]
\label{def:D4}
The \emph{$D_4$ root system} is
\[
  \Phi(D_4) := \{\pm e_i \pm e_j : 1 \leq i < j \leq 4\} \subset U,
\]
consisting of $4 \binom{4}{2} = 24$ vectors of squared length $2$
in $U \cong \R^4$.
\end{definition}

\subsection{The inclusion theorem}

\begin{proposition}[Strict $D_4 \subset E_8$ inclusion]
\label{prop:D4-inclusion}
With the normalizations of \S\ref{sec:E8}:
\begin{enumerate}
\item $\Phi(D_4) = \Phi(E_8) \cap U$.
\item $\Phi(D_4)$ is a reduced, irreducible crystallographic root
  system of rank $4$ and order $24$.
\item $\Phi(D_4) \subsetneq \Phi(E_8)$, i.e.\ the inclusion is
  strict.
\end{enumerate}
\end{proposition}

\begin{proof}
\emph{(1) $\Phi(D_4) \subseteq \Phi(E_8) \cap U$.} Every element
$\pm e_i \pm e_j$ with $1 \leq i < j \leq 4$ has zero components in
positions $5, \ldots, 8$, so it lies in $U$. It is also a
$D_8$-type root of $\Phi(E_8)$ by Definition~\ref{def:E8}. Hence
$\pm e_i \pm e_j \in \Phi(E_8) \cap U$.

\emph{$\Phi(E_8) \cap U \subseteq \Phi(D_4)$.} Let $r \in \Phi(E_8)
\cap U$. Then $r$ has zero components in positions $5, \ldots, 8$
(by membership in $U$). We consider the two types of roots in
$\Phi(E_8)$:
\begin{itemize}
\item \textbf{$D_8$-type.} $r = \pm e_i \pm e_j$ with $1 \leq i <
  j \leq 8$. For $r \in U$, the nonzero components of $r$ must be
  in positions $1, \ldots, 4$, so $1 \leq i < j \leq 4$. Hence
  $r \in \Phi(D_4)$.
\item \textbf{Half-integer.} $r = \frac{1}{2}(\epsilon_1, \ldots,
  \epsilon_8)$ with $\epsilon_i \in \{+1, -1\}$ and an even number
  of minus signs. Every component of $r$ has absolute value
  $\frac{1}{2} \neq 0$; in particular, the components in positions
  $5, \ldots, 8$ are non-zero, so $r \notin U$.
\end{itemize}
So $\Phi(E_8) \cap U$ contains no half-integer roots, and the
$D_8$-type roots in $U$ are exactly $\Phi(D_4)$. Hence
$\Phi(E_8) \cap U \subseteq \Phi(D_4)$.

\emph{(2)} $\Phi(D_4)$ is the root system of type $D_4$ in the
standard model \cite[\S 12.1]{Humphreys}: rank $4$, order $24$,
simple roots $\{e_1 - e_2, e_2 - e_3, e_3 - e_4, e_3 + e_4\}$.
Irreducibility and crystallographicity are classical.

\emph{(3)} $|\Phi(E_8)| = 240 > 24 = |\Phi(D_4)|$, so the
inclusion is strict.
\end{proof}

\begin{remark}[Triality]
\label{rmk:triality}
The Weyl group $W(D_4)$ has the outer automorphism group $S_3$
(triality), permuting the three irreducible $8$-dimensional real
representations of $\Spin(8)$ \cite[\S 20.3]{FultonHarris}. This is
classical background; triality plays a role in \S\ref{sec:8-sector}
in linking the octonionic $8$-sector to the $D_4$ picture, but P2
does not claim any cascade-specific use of triality beyond
compatibility with the standard $\Spin(8)$-representation theory.
\end{remark}

%=====================================================================
\section{Icosians and the $E_8 \to H_4$ Projection}
\label{sec:H4}
%=====================================================================

This is the delicate section. The key technical point is that
$H_4$ is non-crystallographic and therefore cannot be a
sub-root-system of any crystallographic root system, in particular
not of $E_8$. What exists is a projection from $E_8$ lattice data
to an $H_4$ module, realized concretely via the icosian ring. We
define this projection explicitly.

\subsection{Quaternions and the icosian ring}

\begin{notation}[Hamilton quaternions]
Let $\HH = \R \oplus \R i \oplus \R j \oplus \R k$ denote the
Hamilton quaternion algebra, with $i^2 = j^2 = k^2 = ijk = -1$.
For $q = a + bi + cj + dk \in \HH$, the conjugate is
$\bar q = a - bi - cj - dk$, the norm is
$N(q) = q \bar q = a^2 + b^2 + c^2 + d^2$, and the trace is
$\operatorname{tr}(q) = q + \bar q = 2a$.
\end{notation}

\begin{definition}[Icosian ring]
\label{def:icosian-ring}
The \emph{icosian ring} $I \subset \HH$ is the subring of $\HH$
generated by the $120$ unit icosians \cite[\S 8.2,
Eq.~(15)]{ConwaySloane}: the unit quaternions
$\{\pm 1, \pm i, \pm j, \pm k\}$; the $16$ vectors
$\tfrac{1}{2}(\pm 1 \pm i \pm j \pm k)$ (all sign combinations);
and the $96$ vectors of the form
$\tfrac{1}{2}(\pm \varphi^{-1} \pm 1 \cdot a_2 \pm \varphi \cdot a_3)$
where $(a_1, a_2, a_3)$ runs over even permutations of
$(1, i, j, k)$ with one entry replaced by $0$ (the precise
enumeration is given in \cite[\S 8.2.2, Table~8.1]{ConwaySloane}).

As a $\Z$-module, $I$ has rank $8$ \cite[Ch.~8,
Eq.~(17)]{ConwaySloane}; a specific $\Z$-basis is exhibited in
\cite[Table~8.1]{ConwaySloane}. Equivalently
\cite[Theorem~4.1]{MoodyPatera}, $I$ has a $\Z$-basis consisting of
eight specific unit icosians; the precise basis depends on a
choice of normalization, and we denote any fixed such basis by
$\{\omega_1, \ldots, \omega_8\} \subset I$. The ring structure on
$I$ is inherited from $\HH$.
\end{definition}

\begin{remark}[On the rationality of $I$]
\label{rmk:I-rationality}
The $\Q$-vector space $I_\Q := I \otimes_\Z \Q$ has $\Q$-dimension
$8$. The icosian ring is \emph{not} of the simple form
$\Z\{1, i, j, k\} \oplus \Z\varphi \cdot \{1, i, j, k\}$: that
$\Z$-module fails to contain the half-integer generators
$\tfrac{1}{2}(\pm 1 \pm i \pm j \pm k)$. The correct $I$ is
strictly larger and includes such half-integer combinations; this
is what makes its trace form unimodular and equal to the $E_8$
form \cite[Ch.~8, Eq.~(17)]{ConwaySloane}. We rely on this
classical structure rather than on any naive coordinatization.
\end{remark}

\begin{definition}[Icosian star-map]
\label{def:icosian-star}
The \emph{icosian star-map} $\ast \colon I \to I$ is the
coefficientwise Galois conjugation
$\varphi \mapsto \bar\varphi = 1 - \varphi$:
\[
  q^\ast := (a_0 + a_0' \bar\varphi) + (a_1 + a_1' \bar\varphi) i
    + (a_2 + a_2' \bar\varphi) j + (a_3 + a_3' \bar\varphi) k.
\]
This is a ring automorphism of $I$ of order $2$, commuting with
quaternionic conjugation $q \mapsto \bar q$.
\end{definition}

\begin{remark}[$\ast$ is the coefficientwise $\sigma$ of P1]
\label{rmk:star-is-sigma}
Under the identification of $I$ with an $8$-dimensional
$\Q$-vector space via the $\Q$-basis
$\{1, i, j, k, \varphi, \varphi i, \varphi j, \varphi k\}$, and the
further identification of $I \otimes_\Z \Q(\varphi)$ with the
scalar extension of a $\Q$-vector space, the star-map $\ast$
coincides with the coefficientwise $\sigma$-action of
\cite[Def.~2.2, Def.~2.5]{SigmaRationality}. We use $\ast$ for the
star-map on $I$ and $\sigma$ for the more general coefficientwise
action per P1; on $I \otimes_\Z \Q(\varphi)$ they agree.
\end{remark}

\subsection{The Euclidean trace form and $E_8$ realization}

\begin{definition}[Euclidean trace form on $I$]
\label{def:trace-form}
The \emph{Euclidean trace form} on $I$ is the $\Z$-bilinear pairing
\[
  \langle p, q \rangle_I := \tfrac{1}{2}\big( \tr(p \bar q)
  + \tr(p^\ast \bar{q}^\ast)\big),
  \qquad p, q \in I.
\]
Equivalently, $\langle p, q \rangle_I = \tr(p\bar q)_{\rm sym}$,
where $\tr(\cdot)_{\rm sym}$ denotes the trace of the symmetric
part under the $\ast$-action (i.e.\ $\tr(\cdot)_{\rm sym} =
\tr \circ \Pi_\sigma$ with $\Pi_\sigma$ of
\cite[Def.~4.1]{SigmaRationality}).
\end{definition}

\begin{fact}[Icosians realize $E_8$; \cite{ConwaySloane, MoodyPatera}]
\label{fact:icosian-E8}
The $\Z$-module $I$ equipped with the Euclidean trace form
$\langle \cdot, \cdot \rangle_I$ is a rank-$8$ Euclidean lattice
isometric to $\Lambda(E_8)$. There exists a $\Z$-module isometry
\[
  \iota_I \colon \Lambda(E_8) \xrightarrow{\;\sim\;} I
\]
carrying the standard $E_8$ form to $\langle \cdot, \cdot \rangle_I$.
The $240$ minimal vectors of $I$ (those $q$ with
$\langle q, q \rangle_I = 2$) correspond bijectively under $\iota_I^{-1}$
to the $240$ roots in $\Phi(E_8)$, and equal the $240$ unit
icosians \cite[Ch.~8, Eq.~(17) and Table~8.1]{ConwaySloane},
\cite[\S 4.1]{MoodyPatera}.
\end{fact}

\begin{notation}[Choice of $\iota_I$]
\label{not:choice-iotaI}
We fix one specific $\iota_I$ from now on, namely the realization
described explicitly in \cite[Ch.~8, Table~8.1]{ConwaySloane}
sending the simple-root basis of $\Lambda(E_8)$ chosen there to a
specified $\Z$-basis of $I$. The downstream constructions of this
paper depend on $\iota_I$ only through its image and through its
$\Z$-linearity; different classical choices of $\iota_I$ differ by
an automorphism of $\Lambda(E_8)$ (an element of $W(E_8)$ acting on
either side), and produce $W(H_4)$-equivalent images on the
$H_4$-side after the projection of \S\ref{subsec:pi-H} below.
\end{notation}

\subsection{The physical-space projection $\pi_H$}
\label{subsec:pi-H}

\begin{definition}[Real embeddings of $\Q(\varphi)$]
\label{def:real-embeddings}
Let $\tau_+, \tau_- \colon \Q(\varphi) \to \R$ denote the two
real embeddings of $\Q(\varphi)$, specified by $\tau_+(\varphi) =
\varphi = (1+\sqrt 5)/2$ and $\tau_-(\varphi) = \bar\varphi =
(1-\sqrt 5)/2$. The pair $(\tau_+, \tau_-)$ is the unique way of
realizing $\Q(\varphi)$ as a subring of $\R \oplus \R$.
\end{definition}

\begin{definition}[Twin map $\eta$]
\label{def:eta}
Every $q \in I$ admits a unique expression as $q = c_0 + c_1 i +
c_2 j + c_3 k$ with $c_0, c_1, c_2, c_3 \in \Q(\varphi)$ inside the
quaternion algebra $\HH \otimes_\R \R$ (since $I \subset \HH$ and
all generators of $I$ have $\Q(\varphi)$-valued components in the
$\{1, i, j, k\}$ basis of $\HH$, by inspection of the generators in
Def.~\ref{def:icosian-ring}). The \emph{twin map}
$\eta \colon I \to \HH \oplus \HH$ is
\[
  \eta(q) := (q_+, q_-),
\]
where $q_\pm := \tau_\pm(c_0) + \tau_\pm(c_1) i + \tau_\pm(c_2) j
+ \tau_\pm(c_3) k$ is obtained by applying $\tau_\pm$ to each
$\Q(\varphi)$-coefficient. The star-map of
Definition~\ref{def:icosian-star} satisfies $\eta(q^\ast) =
(q_-, q_+)$, i.e.\ $\eta$ interchanges the two $\HH$-factors under
the star-map; equivalently $\operatorname{swap} \circ \eta = \eta
\circ \ast$, where $\operatorname{swap}(p, q) := (q, p)$.
\end{definition}

\begin{definition}[$H_4$ root system]
\label{def:H4}
We write $\Phi(H_4) \subset \HH \cong \R^4$ for the standard
coordinate realization of the $H_4$ root system: the $120$ unit
icosians of $\HH$, equivalently the $120$ vertices of the $600$-cell
in the standard coordinates \cite[\S 22.7]{CoxeterRegularPolytopes},
\cite[\S 4.1]{MoodyPatera}. We refer to this set as the
\emph{$120$ unit icosians} or as $\Phi(H_4)$ interchangeably; both
mean the same $120$-element subset of $\HH$.
\end{definition}

\begin{definition}[Physical-space projection]
\label{def:pr-phys}
The \emph{physical-space projection}
$\operatorname{pr}_{\rm phys} \colon \HH \oplus \HH \to \HH$ is
$\operatorname{pr}_{\rm phys}(q_+, q_-) := q_+$, i.e.\ projection
onto the first factor.
\end{definition}

\begin{definition}[The $E_8 \to H_4$ projection]
\label{def:pi-H}
The \emph{$E_8 \to H_4$ projection} is the composition
\[
  \pi_H \colon \Lambda(E_8) \xrightarrow{\;\iota_I\;} I
  \xrightarrow{\;\eta\;} \HH \oplus \HH
  \xrightarrow{\;\operatorname{pr}_{\rm phys}\;} \HH,
  \qquad \pi_H(x) := (\iota_I(x))_+.
\]
The image $\pi_H(\Lambda(E_8)) \subset \HH \cong \R^4$ is the set
$\{q_+ : q \in I\}$ obtained by applying $\tau_+$ coefficientwise
to every icosian.
\end{definition}

\begin{theorem}[$E_8 \to H_4$ projection: structural properties]
\label{thm:pi-H}
The map $\pi_H \colon \Lambda(E_8) \to \HH \cong \R^4$ of
Definition~\ref{def:pi-H} satisfies:
\begin{enumerate}
\item $\pi_H$ is $\Z$-linear.
\item $\pi_H(\Lambda(E_8))$ is the additive group
  $\{q_+ : q \in I\}$, which is a dense subgroup of $\HH$.
\item The image of the $E_8$ minimal shell $\{r \in \Phi(E_8) :
  \langle r, r \rangle = 2\}$ under $\pi_H$ is exactly the
  $120$-element set $\Phi(H_4)$ of Definition~\ref{def:H4}.
\end{enumerate}
$H_4$ is non-crystallographic, in the sense of
Lemma~\ref{lem:H4-noncrystallographic} below; in particular
$\Phi(H_4)$ is \emph{not} a sub-root-system of $\Phi(E_8)$.
\end{theorem}

\begin{proof}
\emph{(1) $\Z$-linearity.} $\iota_I$ is $\Z$-linear by
Fact~\ref{fact:icosian-E8}. $\eta$ is $\Z$-linear because $\tau_\pm
\colon \Q(\varphi) \to \R$ are $\Q$-linear ring homomorphisms,
and $\eta$ acts coefficientwise on the $\Z[\varphi]$-coefficients
of icosians. $\operatorname{pr}_{\rm phys}$ is the projection onto
the first $\HH$-factor of $\HH \oplus \HH$, hence $\Z$-linear.
Composition: $\pi_H = \operatorname{pr}_{\rm phys} \circ \eta \circ
\iota_I$ is $\Z$-linear.

\emph{(2) Image and density.} By Definition~\ref{def:pi-H},
$\pi_H(\Lambda(E_8)) = \{q_+ : q \in I\}$ since $\iota_I$ is
surjective onto $I$. We show this set is dense in $\HH \cong \R^4$.

By \cite[Theorem~4.1 and Proposition~4.2]{MoodyPatera}, the
$\tau_+$-image of $I$ in $\HH$ is precisely the
$\Z[\varphi]$-module generated by $\{1, i, j, k\}$, where
$\varphi = (1+\sqrt 5)/2 \in \R$ now denotes the real number:
\[
  \pi_H(\Lambda(E_8)) = \{q_+ : q \in I\} =
    \Z[\varphi]\{1, i, j, k\} \subset \HH.
\]
This is the classical ``cut-and-project'' image of the icosian
ring under $\tau_+$ \cite[\S 4.1]{MoodyPatera}.

The $\Z[\varphi]$-module $\Z[\varphi]\{1, i, j, k\}$ has
$\Z$-rank $8$ in $\HH \cong \R^4$ (rank $2$ in each of the four
coordinate directions, since $\Z[\varphi] = \Z + \Z\varphi$ has
$\Z$-rank $2$). Since $\varphi$ is irrational
($\varphi^2 = \varphi + 1$ has discriminant $5$, not a square in
$\Q$), the additive subgroup $\Z + \Z\varphi \subset \R$ is dense
in $\R$ \cite[Theorem~A.1]{MoodyPatera} (an instance of the
classical result that $\Z + \Z\alpha$ is dense in $\R$ for any
irrational $\alpha$). Density in each of the four coordinates of
$\HH$ implies density of the product $\Z[\varphi]\{1, i, j, k\}$
in $\HH \cong \R^4$. Hence $\pi_H(\Lambda(E_8))$ is dense in
$\HH$.

\emph{(3) Image of the minimal shell.} The $E_8$ minimal shell
under $\iota_I$ corresponds to the $240$ minimal vectors of $I$,
which by Fact~\ref{fact:icosian-E8} are the $240$ unit icosians.
The $240$ unit icosians come in $120$ pairs $(q, q^\ast)$ each
mapping under $\eta$ to the swap pair $((q_+, q_-), (q_-, q_+))
\in \HH \oplus \HH$ (with $q^\ast \neq q$ when $q \notin
\Z[\varphi]$). Applying $\operatorname{pr}_{\rm phys}$ to each
pair gives the same first-coordinate set: the $120$-element
projection of these pairs is exactly the set of $120$ unit
icosians in $\HH$, which by Definition~\ref{def:H4} is
$\Phi(H_4)$. (For the $\Z[\varphi]$-rational unit icosians ---
i.e.\ those fixed by $\ast$ --- the pair degenerates and the count
is preserved via the standard $\tau_+$-image \cite[\S 4.1,
Eq.~(4.5)]{MoodyPatera}.)
\end{proof}

\begin{lemma}[$H_4$ is non-crystallographic]
\label{lem:H4-noncrystallographic}
There is no $W(H_4)$-stable lattice $L \subset \R^4$. In particular,
$\Phi(H_4)$ is not a sub-root-system of any crystallographic root
system of rank $4$ in $\R^4$, including any embedded $\Phi(D_4) \subset
\R^4$.
\end{lemma}

\begin{proof}
$W(H_4)$ contains rotations of order $5$ (it is the symmetry group
of the regular pentagon-faced $120$-cell, with $5$-fold faces).
Crystallographic restrictions in $\R^4$ \cite[Theorem~6.4]{Humphreys}
allow rotational orders only in $\{1, 2, 3, 4, 6\}$ for any
lattice-preserving isometry of a rank-$4$ lattice; order $5$ is
excluded. Hence $W(H_4)$ cannot preserve any rank-$4$ lattice.
\end{proof}

\begin{remark}[On part (2) and crystallographic embeddings]
\label{rmk:density-noembedding}
Combining (2) and Lemma~\ref{lem:H4-noncrystallographic} clarifies
the relationship between $E_8$ and $H_4$: the additive image
$\pi_H(\Lambda(E_8))$ is dense in $\R^4$, which is consistent with
the projective fact that no crystallographic lattice in $\R^4$
contains $\Phi(H_4)$. In particular, the $E_8$-to-$H_4$
correspondence cannot be presented as a sub-root-system inclusion
even after rescaling. Theorem~\ref{thm:pi-H}(3) is the precise
statement of what \emph{is} true: the minimal shell maps onto
$\Phi(H_4)$ as sets under the projection.
\end{remark}

\begin{remark}[Why this is a projection, not an inclusion]
\label{rmk:projection-not-inclusion}
Theorem~\ref{thm:pi-H} expresses the $E_8$-$H_4$ relationship as a
projection from an $8$-dimensional crystallographic object
($\Lambda(E_8)$) to a $4$-dimensional non-crystallographic object
($H_4$). This is often described informally as ``$H_4 \subset
E_8$''; that description is algebraically misleading. The precise
statement is Theorem~\ref{thm:pi-H}(3): $\Phi(H_4)$ equals the
image of the minimal shell under $\pi_H$, not a sub-root-system of
$\Phi(E_8)$.
\end{remark}

%=====================================================================
\section{The $24$-Cell Inside the $600$-Cell}
\label{sec:24-600}
%=====================================================================

This section isolates an explicit $24$-vertex subset of the
$600$-cell that is a regular $24$-cell, isometric to the $D_4$ root
polytope. The proof is by direct coordinate inspection.

\subsection{Standard $600$-cell coordinates}

\begin{definition}[$V_{600}$, the $600$-cell vertex set]
\label{def:V600}
Working in $\R^4$ with coordinates $(x_1, x_2, x_3, x_4)$ and
$\varphi = (1+\sqrt 5)/2$, define $V_{600}$ as the union of the
following three classes of vectors
\cite[\S 22.7]{CoxeterRegularPolytopes}, \cite[\S 3.1]{MoodyPatera}:
\begin{enumerate}
\item \textbf{Axis points ($8$):}
  all permutations of $(\pm 1, 0, 0, 0)$ (i.e.\ the eight vectors
  $\pm e_1, \pm e_2, \pm e_3, \pm e_4$).
\item \textbf{Half-type points ($16$):}
  $\frac{1}{2}(\pm 1, \pm 1, \pm 1, \pm 1)$, with independent sign
  choices in each of the four coordinates; $2^4 = 16$ vectors.
\item \textbf{Golden points ($96$):}
  all even permutations of
  $\tfrac{1}{2}(0, \pm 1, \pm \varphi^{-1}, \pm \varphi)$, with
  independent signs on the three nonzero coordinates. There are
  $\tfrac{4!}{2} = 12$ even permutations of the four positions and
  $2^3 = 8$ independent sign choices on the nonzero coordinates,
  giving $12 \cdot 8 = 96$ vectors.
\end{enumerate}
The total is $|V_{600}| = 8 + 16 + 96 = 120$, matching the standard
vertex count of the $600$-cell \cite[\S 22.7]{CoxeterRegularPolytopes}.
\end{definition}

\begin{remark}[Normalization]
With the coordinates above, every $v \in V_{600}$ satisfies
$\|v\|^2 = 1$, i.e.\ $V_{600}$ lies on the unit $3$-sphere in $\R^4$.
\end{remark}

\subsection{The inscribed $24$-cell}

\begin{definition}[$V_{24}$, the inscribed $24$-cell vertex set]
\label{def:V24}
Let
\[
  V_{24} := \big\{\text{axis points}\big\}
    \cup \big\{\text{half-type points}\big\} \subset V_{600}.
\]
Explicitly: $V_{24}$ consists of the $8$ vectors of the form
$(\pm 1, 0, 0, 0)$ and permutations, together with the $16$
vectors $\frac{1}{2}(\pm 1, \pm 1, \pm 1, \pm 1)$ (all signs).
$|V_{24}| = 24$.
\end{definition}

\begin{proposition}[$24$-cell inside $600$-cell]
\label{prop:24-in-600}
The rescaled $D_4$ root system $\tilde\Phi(D_4) := \{r/\sqrt 2 :
r \in \Phi(D_4)\}$ is a $24$-element set of unit vectors in $\R^4$.
There exists an explicit orthogonal transformation $R \in
\mathrm{O}(4)$ such that $R\big(\tilde\Phi(D_4)\big) = V_{24}$. In
particular $V_{24}$ and $\tilde\Phi(D_4)$ are congruent.
\end{proposition}

\begin{proof}
Define the block-orthogonal matrix
\[
  R := \frac{1}{\sqrt 2}
  \begin{pmatrix}
    1 & 1 & 0 & 0 \\
    1 & -1 & 0 & 0 \\
    0 & 0 & 1 & 1 \\
    0 & 0 & 1 & -1
  \end{pmatrix}.
\]
Direct computation gives $R^T R = I_4$, so $R \in \mathrm{O}(4)$.

We verify $R(\tilde\Phi(D_4)) = V_{24}$ by direct evaluation on the
$24$ unit vectors
$\tilde\Phi(D_4) = \big\{\tfrac{1}{\sqrt 2}(\pm e_i \pm e_j) :
1 \leq i < j \leq 4 \big\}$. The six unordered index pairs
$\{i, j\}$ split into two classes:

\emph{Class A: $\{i, j\} \in \{\{1, 2\}, \{3, 4\}\}$
($R$-block-internal pairs).} Direct computation:
\begin{align*}
  R\big(\tfrac{1}{\sqrt 2}(e_1 + e_2)\big) &=
    \tfrac{1}{\sqrt 2} \cdot \tfrac{1}{\sqrt 2}(2, 0, 0, 0) =
    e_1, \\
  R\big(\tfrac{1}{\sqrt 2}(e_1 - e_2)\big) &= e_2, \\
  R\big(\tfrac{1}{\sqrt 2}(e_3 + e_4)\big) &= e_3, \\
  R\big(\tfrac{1}{\sqrt 2}(e_3 - e_4)\big) &= e_4.
\end{align*}
Including the $\pm$ signs, the $8$ rescaled roots in Class A map
bijectively to the $8$ axis points $\{\pm e_1, \pm e_2, \pm e_3,
\pm e_4\}$ of $V_{24}$.

\emph{Class B: $\{i, j\} \in \{\{1,3\}, \{1,4\}, \{2,3\}, \{2,4\}\}$
(cross-block pairs).} Direct computation:
\begin{align*}
  R\big(\tfrac{1}{\sqrt 2}(e_1 + e_3)\big) &=
    \tfrac{1}{2}(1, 1, 1, 1), &
  R\big(\tfrac{1}{\sqrt 2}(e_1 + e_4)\big) &=
    \tfrac{1}{2}(1, 1, 1, -1), \\
  R\big(\tfrac{1}{\sqrt 2}(e_1 - e_3)\big) &=
    \tfrac{1}{2}(1, 1, -1, -1), &
  R\big(\tfrac{1}{\sqrt 2}(e_1 - e_4)\big) &=
    \tfrac{1}{2}(1, 1, -1, 1), \\
  R\big(\tfrac{1}{\sqrt 2}(e_2 + e_3)\big) &=
    \tfrac{1}{2}(1, -1, 1, 1), &
  R\big(\tfrac{1}{\sqrt 2}(e_2 + e_4)\big) &=
    \tfrac{1}{2}(1, -1, 1, -1), \\
  R\big(\tfrac{1}{\sqrt 2}(e_2 - e_3)\big) &=
    \tfrac{1}{2}(1, -1, -1, -1), &
  R\big(\tfrac{1}{\sqrt 2}(e_2 - e_4)\big) &=
    \tfrac{1}{2}(1, -1, -1, 1).
\end{align*}
The $8$ images all have first coordinate $+\tfrac{1}{2}$, and their
sign patterns on the remaining three coordinates realize all $2^3
= 8$ combinations. Including the negatives $-r$ for each rescaled
root $r$ (which produce images with first coordinate
$-\tfrac{1}{2}$ and exhaust the remaining sign patterns), the $16$
rescaled roots in Class B map bijectively to the $16$ half-type
vertices $\tfrac{1}{2}(\pm 1, \pm 1, \pm 1, \pm 1)$ of $V_{24}$.

Classes A and B together cover $\tilde\Phi(D_4)$ ($8 + 16 = 24$
vectors), and their images cover $V_{24}$ ($8$ axis $+\, 16$
half-type $= 24$ vertices). Hence $R(\tilde\Phi(D_4)) = V_{24}$
as sets, and since $R \in \mathrm{O}(4)$ preserves Euclidean
distances, $V_{24}$ and $\tilde\Phi(D_4)$ are congruent.
\end{proof}

\begin{corollary}[$V_{24}$ is a regular $24$-cell]
\label{cor:V24-24cell}
$V_{24}$ is the vertex set of a regular $24$-cell in $\R^4$.
\end{corollary}

\begin{proof}
$\tilde\Phi(D_4)$ is the standard unit-radius realization of the
$D_4$ root polytope. The $D_4$ root polytope is the regular
$24$-cell with $24$ vertices on the unit sphere
\cite[\S 22.7]{CoxeterRegularPolytopes} (which gives the same
$24$-cell by an alternative coordinate model; the two models are
related by the rotation $R$ of
Proposition~\ref{prop:24-in-600}). By
Proposition~\ref{prop:24-in-600}, $V_{24}$ is the image of
$\tilde\Phi(D_4)$ under an orthogonal transformation of $\R^4$,
hence itself a regular $24$-cell.
\end{proof}

\begin{remark}[Why this is an inscription, not a quotient]
\label{rmk:inscription-not-quotient}
Proposition~\ref{prop:24-in-600} says $V_{24}$ is a subset of
$V_{600}$, not that $V_{600}$ is a fibre bundle over $V_{24}$ with
``fibres of size $5$''. The $2I/2T$ coset language of earlier
cascade notes describes a related combinatorial structure on the
binary icosahedral and tetrahedral groups inside $\Sp(1) \subset
\HH$, but that language is not needed here and is not used as a
theorem statement. The paper uses only
Proposition~\ref{prop:24-in-600}: $V_{24} \subset V_{600}$,
isometric to the $24$-cell, $24$ vertices.
\end{remark}

%=====================================================================
\section{The $600$-Cell Spectral and Shell Structure}
\label{sec:600cell-spectral}
%=====================================================================

The preceding sections constructed the $600$-cell vertex set
$V_{600}$ as a finite subset of the unit $3$-sphere via explicit
coordinates (Definition~\ref{def:V600}), and exhibited an inscribed
$24$-cell $V_{24} \subset V_{600}$. This section collects the
spectral, shell, and representation-theoretic data on the
$600$-cell vertex graph that downstream papers draw on (Papers V
and XXII of the cascade-correspondence programme, the
cascade-refinement-spaces paper P3 for its Markov-chain
substrate, and the Foundations paper for its spectral
decomposition). Every numerical claim in this section is either an
exact finite-arithmetic verification or an explicit import from a
classical source that is cited at the point of use.

\subsection{Icosian group structure under Hamilton multiplication}

\begin{definition}[Hamilton multiplication on $V_{600}$]
\label{def:ham-600}
Regard each vertex $v = (v_1, v_2, v_3, v_4) \in V_{600} \subset
\R^4$ as a unit quaternion $v_1 + v_2 i + v_3 j + v_4 k \in
\HH$. Hamilton multiplication $(v, w) \mapsto v \cdot w$ then
takes two unit quaternions to a unit quaternion. The vertex
identification with unit icosians used here is the standard one,
as in \cite[\S 3.1]{MoodyPatera} and
\cite[Ch.~8]{ConwaySloane}; each coordinate lies in
$\Q(\sqrt{5}) \subset \Q(\varphi)$, so Hamilton products remain in
the same ring.
\end{definition}

\begin{theorem}[Icosian group closure and identification with $2I$,
computer-assisted]
\label{thm:icosian-closure}
The vertex set $V_{600}$ is closed under Hamilton multiplication,
and, as a subgroup of $\mathrm{Sp}(1) \subset \HH$, is isomorphic
to the binary icosahedral group $2I$ of order $120$.
\end{theorem}

\begin{proof}
\emph{(L1) Exact vertex construction.} Build each of the $120$
vertices of Definition~\ref{def:V600} with every coordinate
represented as an exact pair $(a, b) \in \Q^2$ denoting $a +
b\sqrt 5$. The script
\verb|papers/paper-xxii/scripts/run_icosian_exact.py| (referenced
here as the canonical P2 implementation, with the same protocol
reused in
\verb|papers/paper-v/scripts/run_rd4_exact_rational.py|) implements
this step.

\emph{(L2) Closure.} Enumerate all $120 \times 120 = 14\,400$
Hamilton products in exact $\Q(\sqrt 5)$-arithmetic and match each
against a vertex by pair equality. The script reports zero
failures: every product is again a vertex. Hence $V_{600}$ is a
finite subgroup of $\mathrm{Sp}(1)$ of order $120$.

\emph{(L3) Classification.} The finite subgroups of
$\mathrm{Sp}(1)$ are classified up to conjugacy in
\cite[Ch.~3]{DuVal1964}: at order $120$, the only finite subgroup
of $\mathrm{Sp}(1)$ is the binary icosahedral group $2I$.
Therefore $V_{600} \cong 2I$ as abstract groups.
\end{proof}

Throughout the remainder of this section we write $V_{600}$ and
$2I$ interchangeably. The step (L2) is computer-assisted in the
sense that a finite enumeration (of $14\,400$ exact pair-of-rational
equalities) is carried out in software; no floating-point
tolerance enters the step. Step (L3) is a purely classical
invocation.

\subsection{Euclidean distance shells and the shell-class bijection}
\label{sec:shells-classes}

\begin{definition}[Euclidean shells from the identity]
\label{def:shells}
For each $v \in V_{600}$, let $\|1 - v\| \in [0, 2]$ be the
Euclidean distance from the identity quaternion $1$ to $v$. The
squared distance satisfies $\|1 - v\|^2 = 2 - 2 v_1 \in
\Q(\sqrt 5)$ with $v_1$ the scalar part of $v$. Group vertices by
exact equality of $\|1 - v\|^2$; this yields the nine \emph{Euclidean
shells} $S_0^{\mathrm{id}}, \ldots, S_8^{\mathrm{id}}$ displayed in
Table~\ref{tab:P2-shells}.
\end{definition}

\begin{table}[ht]
\centering
\caption{The nine Euclidean distance shells of $V_{600}$ from the
identity, listed by exact squared distance, Euclidean distance,
spherical angle, and shell size. Reproduced in exact
$\Q(\sqrt 5)$-arithmetic by
\texttt{run\_icosian\_exact.py}.}
\label{tab:P2-shells}
\begin{tabular}{@{}cllllr@{}}
\toprule
$k$ & Squared distance & Distance & Angle & $\cos(\text{angle})$ & $|S_k^{\mathrm{id}}|$ \\
\midrule
$0$ & $0$                  & $0$              & $0^\circ$   & $1$              & $1$ \\
$1$ & $(3-\sqrt 5)/2 = 1/\varphi^2$ & $1/\varphi$     & $36^\circ$  & $\varphi/2$      & $12$ \\
$2$ & $1$                  & $1$              & $60^\circ$  & $1/2$            & $20$ \\
$3$ & $(5-\sqrt 5)/2 = 3-\varphi$ & $\sqrt{3-\varphi}$ & $72^\circ$  & $1/(2\varphi)$ & $12$ \\
$4$ & $2$                  & $\sqrt 2$        & $90^\circ$  & $0$              & $30$ \\
$5$ & $(3+\sqrt 5)/2 = \varphi^2$ & $\varphi$       & $108^\circ$ & $-1/(2\varphi)$  & $12$ \\
$6$ & $3$                  & $\sqrt 3$        & $120^\circ$ & $-1/2$           & $20$ \\
$7$ & $(5+\sqrt 5)/2 = 2+\varphi$ & $\sqrt{3+\varphi}$ & $144^\circ$ & $-\varphi/2$ & $12$ \\
$8$ & $4$                  & $2$              & $180^\circ$ & $-1$             & $1$ \\
\bottomrule
\end{tabular}
\end{table}

\begin{theorem}[Shell = conjugacy class, computer-assisted]
\label{thm:shell-class}
The nine Euclidean shells $S_k^{\mathrm{id}}$ are in size-preserving
bijection with the nine conjugacy classes of $2I$. Explicitly: for
each $k$, the set $S_k^{\mathrm{id}} \subset V_{600}$ is a single
$2I$-conjugacy class, and the shell-size multiset $\{1, 12, 20,
12, 30, 12, 20, 12, 1\}$ matches the conjugacy-class-size multiset
of $2I$.
\end{theorem}

\begin{proof}
\emph{(L3a) Conjugacy classes of $2I$.} Using the exact $\Q(\sqrt
5)$-arithmetic of Theorem~\ref{thm:icosian-closure}, enumerate for
each $g \in V_{600}$ the orbit $\{h g h^{-1} : h \in V_{600}\}$
under Hamilton conjugation, and bucket vertices by exact
coordinate equality. The script
\texttt{run\_icosian\_exact.py} reports nine disjoint classes of
sizes $\{1, 1, 12, 12, 12, 12, 20, 20, 30\}$, summing to $120$
(the standard class structure of $2I$).

\emph{(L3b) Euclidean shell decomposition.} For each $v \in
V_{600}$ compute $\|1 - v\|^2 = 2 - 2 v_1 \in \Q(\sqrt 5)$ and
bucket by exact pair equality. Output: nine shells with squared
distances and sizes as in Table~\ref{tab:P2-shells}; sizes $\{1,
12, 20, 12, 30, 12, 20, 12, 1\}$ summing to $120$.

\emph{(L3c) Shell-class matching.} The shell-size multiset matches
the class-size multiset. A direct exact comparison of the nine
shell-index-sets against the nine conjugacy-class-index-sets in
the same script shows equality pair-by-pair: each Euclidean
distance shell from the identity coincides with exactly one
$2I$-conjugacy class.

\emph{Conceptual reason.} The inner product on $\HH$ satisfies
$\|1 - v\|^2 = 2 - 2\,\mathrm{Re}(v) = 2 - 2\,\mathrm{Re}(h v
h^{-1})$ for any unit quaternion $h$ (since
$\mathrm{Re}$ is conjugation-invariant on $\HH$). Therefore every
$2I$-conjugacy class lies in a single Euclidean shell from $1$.
The finite enumeration in (L3c) confirms that each shell contains
exactly one class, i.e.\ no shell merges two distinct classes.
\end{proof}

\begin{corollary}[Shell-adjacency class sums are central]
\label{cor:class-sums-central}
For each $k \in \{1, \ldots, 8\}$, the \emph{shell-$k$ class sum}
\[
  C_k \;:=\; \sum_{s \in S_k^{\mathrm{id}}} s \;\in\; \C[2I]
\]
is a central element of the group algebra $\C[2I]$. Consequently,
the \emph{shell-$k$ adjacency matrix} $\mathbf{A}_k$ on
$\C^{V_{600}}$ defined by $(\mathbf{A}_k)_{v, w} = 1$ if
$\|v - w\| = d_k$ and $0$ otherwise acts, in the identification
$\C^{V_{600}} \cong \C[2I]$ afforded by the left-regular action,
as \emph{right-convolution by $C_k$}; and by Schur's lemma it acts
by a scalar on each isotypic component of $\C[2I]$.
\end{corollary}

\begin{proof}
For unit quaternions $v, w$, $\|v - w\|^2 = 2 - 2 \mathrm{Re}(v
w^{-1})$, so $\|v - w\| = d_k$ iff $v w^{-1} \in S_k^{\mathrm{id}}$.
Hence $(\mathbf{A}_k)_{v, w} = \mathbf{1}[v w^{-1} \in
S_k^{\mathrm{id}}] = \mathbf{1}[v w^{-1} \in C_k]$ in class-sum
notation. In the left-regular representation, this is precisely
right-convolution of $\delta_w$ by $C_k$. Each $S_k^{\mathrm{id}}$
is a conjugacy class (Theorem~\ref{thm:shell-class}), so $C_k$ is
central; Schur's lemma then gives the scalar action on each
isotypic component of $\C[2I] = \bigoplus_\rho (\dim\rho) \cdot
V_\rho$.
\end{proof}

\subsection{Laplacian spectrum and the $94 + 26$ integer /
irrational decomposition}
\label{sec:spectrum-P2}

\begin{fact}[Laplacian spectrum of the $600$-cell vertex graph,
imported computation]
\label{prop:spectrum-P2}
Let $\mathbf{A}_1$ be the nearest-neighbour adjacency matrix of
the $600$-cell vertex graph (shell $k = 1$, degree $12$), and let
$L = 12 I - \mathbf{A}_1$ be the combinatorial graph Laplacian.
Then $L$ has the nine distinct eigenvalues and multiplicities of
Table~\ref{tab:P2-spectrum}.

This fact is imported. Its two-step derivation is standard and
not reproduced internally in this paper: (i) by
Corollary~\ref{cor:class-sums-central}, $\mathbf{A}_1$ acts on
each isotypic component $(\dim \rho) \cdot V_\rho$ of $\C[2I]$ by
the scalar $\alpha_\rho = (\dim \rho)^{-1} \sum_{s \in
S_1^{\mathrm{id}}} \chi_\rho(s)$; (ii) evaluating $\alpha_\rho$
for each of the nine irreducible representations of $2I$ requires
the specific entries of the standard $2I$ character table, which
we import from \cite[Ch.~3]{DuVal1964} (Du~Val's classification
of the binary polyhedral groups already contains their character
tables). The numerical values of Table~\ref{tab:P2-spectrum} are
independently cross-checked at floating-point precision by
\verb|papers/paper-xxii/scripts/run_600cell_spectral_verification.py|
and by the diagonalisation carried out in the spectral-check
branch of \texttt{run\_icosian\_exact.py}. No part of this
character-table step is re-proved from first principles in P2.
\end{fact}

\begin{table}[ht]
\centering
\caption{Complete Laplacian spectrum of the $600$-cell vertex
graph. Multiplicities are $(\dim\rho)^2$ for the nine irreducible
representations of $2I$.}
\label{tab:P2-spectrum}
\begin{tabular}{@{}clllrl@{}}
\toprule
Index & $2I$ irrep $\rho$ & $\dim \rho$ & Laplacian eigenvalue $\lambda$ & Multiplicity & Integer? \\
\midrule
$0$ & trivial       & $1$ & $0$              & $1  = 1^2$ & Yes \\
$1$ & $\mathbf{2}$  & $2$ & $9 - 3\sqrt 5$   & $4  = 2^2$ & No  \\
$2$ & $\mathbf{3}$  & $3$ & $10 - 2\sqrt 5$  & $9  = 3^2$ & No  \\
$3$ & $\mathbf{4}$  & $4$ & $9$              & $16 = 4^2$ & Yes \\
$4$ & $\mathbf{5}$  & $5$ & $12$             & $25 = 5^2$ & Yes \\
$5$ & $\mathbf{6}$  & $6$ & $14$             & $36 = 6^2$ & Yes \\
$6$ & $\mathbf{3'}$ & $3$ & $10 + 2\sqrt 5$  & $9  = 3^2$ & No  \\
$7$ & $\mathbf{4'}$ & $4$ & $15$             & $16 = 4^2$ & Yes \\
$8$ & $\mathbf{2'}$ & $2$ & $9 + 3\sqrt 5$   & $4  = 2^2$ & No  \\
\bottomrule
\end{tabular}
\end{table}

\begin{corollary}[Integer / irrational isotypic decomposition]
\label{cor:94-26-split}
Let $W_{\mathrm{int}} \subset \C^{V_{600}}$ be the direct sum of
those isotypic components whose Laplacian eigenvalue in
Table~\ref{tab:P2-spectrum} is an integer, and let $W_{\mathrm{irr}}
\subset \C^{V_{600}}$ be the direct sum of those whose Laplacian
eigenvalue is irrational (i.e.\ lies in $\Q(\sqrt 5) \setminus
\Q$). Then
\begin{equation}
  \C^{V_{600}} \;=\; W_{\mathrm{int}} \oplus W_{\mathrm{irr}},
  \qquad
  \dim W_{\mathrm{int}} \;=\; 1 + 16 + 25 + 36 + 16 \;=\; 94,
  \qquad
  \dim W_{\mathrm{irr}} \;=\; 4 + 9 + 9 + 4 \;=\; 26,
\end{equation}
and for every polynomial $p$ in the nine shell-adjacency matrices
$\mathbf{A}_1, \ldots, \mathbf{A}_8$,
\begin{equation}
  P_{\mathrm{int}} \, p(\mathbf{A}_1, \ldots, \mathbf{A}_8) \,
  P_{\mathrm{irr}} \;=\; 0,
\end{equation}
where $P_{\mathrm{int}}, P_{\mathrm{irr}}$ are the orthogonal
projectors onto $W_{\mathrm{int}}$ and $W_{\mathrm{irr}}$.
\end{corollary}

\begin{proof}
Each $\mathbf{A}_k$ acts as a scalar on every isotypic component
(Corollary~\ref{cor:class-sums-central}); both $W_{\mathrm{int}}$
and $W_{\mathrm{irr}}$ are unions of isotypic components; hence
each $\mathbf{A}_k$, and every polynomial in them, preserves both
summands. The dimension count is read off
Table~\ref{tab:P2-spectrum}.
\end{proof}

\begin{remark}[Provenance of the $94{+}26$ split]
The invariance of $W_{\mathrm{int}}$ and $W_{\mathrm{irr}}$ under
the shell-adjacency algebra is proved here without reference to
the $2I$ character table: it follows from
Theorem~\ref{thm:shell-class} (shells are conjugacy classes) and
Corollary~\ref{cor:class-sums-central} (class sums are central).
What the character-table import supplies is the \emph{labelling}
of individual isotypic components as integer- or
irrational-eigenvalue. Downstream papers that merely need the
algebra to preserve some $94{+}26$ split therefore depend only on
the $\Q(\sqrt 5)$-exact arithmetic step; downstream papers that
need the specific integer and irrational eigenvalues $\{0, 9, 12,
14, 15\}$ and $\{9 \pm 3\sqrt 5, 10 \pm 2\sqrt 5\}$ inherit the
classical-character-table dependence recorded in the exclusion
list of \S\ref{sec:intro}.
\end{remark}

\subsection{BFS shell structure and dual subpopulations}
\label{sec:BFS-P2}

\begin{proposition}[BFS shells of the $600$-cell vertex graph]
\label{prop:BFS-shells}
Fix any vertex $v_0 \in V_{600}$ as a source. The BFS (graph
distance) partition from $v_0$ gives six shells
$B_0, B_1, \ldots, B_5$ with sizes
\[
  |B_0|, |B_1|, |B_2|, |B_3|, |B_4|, |B_5|
  \;=\; 1,\;12,\;32,\;42,\;32,\;1.
\]
The shell sizes are source-independent (by the $H_4$
vertex-transitivity of the $600$-cell, equivalently the
$2I$-homogeneity under left-multiplication), and they refine the
Euclidean shells of Table~\ref{tab:P2-shells} as follows: $B_0 =
S_0^{\mathrm{id}}$, $B_1 = S_1^{\mathrm{id}}$, $B_2 =
S_2^{\mathrm{id}} \cup S_3^{\mathrm{id}}$, $B_3 = S_4^{\mathrm{id}}
\cup S_5^{\mathrm{id}}$, $B_4 = S_6^{\mathrm{id}} \cup
S_7^{\mathrm{id}}$, $B_5 = S_8^{\mathrm{id}}$.
\end{proposition}

\begin{proof}
Direct breadth-first search on the $120$-vertex graph starting
from any $v_0$, carried out in
\verb|papers/paper-v/scripts/run_rigorous_verification.py|. The
BFS-to-Euclidean-shell refinement is read off by comparing BFS
distances with the exact Euclidean shell assignment of
Theorem~\ref{thm:shell-class}.
\end{proof}

\begin{definition}[Icosahedral and dodecahedral subpopulations]
\label{def:ico-dod}
For each shell $B_d \in \{B_2, B_3, B_4\}$, partition $B_d$ by the
triple of integer \emph{intersection numbers} $(a, b, c) :=
(|N(v) \cap B_d|,\, |N(v) \cap B_{d+1}|,\, |N(v) \cap B_{d-1}|)$
for $v \in B_d$, where $N(v)$ is the nearest-neighbour set. Call
the vertices with $|N(v) \cap B_d| = 5$ \emph{icosahedral} and
the vertices with $|N(v) \cap B_d| = 6$ \emph{dodecahedral}
within shell $B_d$.
\end{definition}

\begin{proposition}[Equitable nine-cell partition with full quotient
matrix]\label{prop:dual-populations}
With the partition of Definition~\ref{def:ico-dod}, the refined
partition
\[
  V_{600} \;=\; B_0 \sqcup B_1 \sqcup B_2^{\mathrm{dod}} \sqcup
    B_2^{\mathrm{ico}} \sqcup B_3^{\mathrm{dod}} \sqcup
    B_3^{\mathrm{ico}} \sqcup B_4^{\mathrm{dod}} \sqcup
    B_4^{\mathrm{ico}} \sqcup B_5
\]
is an equitable $9$-cell partition of the $600$-cell vertex graph
in the sense of \cite[\S 9.3]{GodsilRoyle2001}: for every pair of
cells $(S, T)$ and every $v \in S$, the integer count of
nearest-neighbour edges from $v$ into $T$ has the same value
$M_{S, T} \in \Z_{\geq 0}$, independent of the choice of $v \in
S$. Specifically, the complete quotient matrix $M = (M_{S,
T})_{S, T}$ is
\begin{equation}\label{eq:quotient-matrix}
M \;=\;
\begin{array}{c|ccccccccc}
 & B_0 & B_1 & B_2^{\mathrm{dod}} & B_2^{\mathrm{ico}} &
   B_3^{\mathrm{dod}} & B_3^{\mathrm{ico}} &
   B_4^{\mathrm{dod}} & B_4^{\mathrm{ico}} & B_5 \\
\hline
B_0                & 0 & 12 & 0 & 0 & 0 & 0 & 0 & 0  & 0 \\
B_1                & 1 & 5  & 5 & 1 & 0 & 0 & 0 & 0  & 0 \\
B_2^{\mathrm{dod}} & 0 & 3  & 3 & 3 & 3 & 0 & 0 & 0  & 0 \\
B_2^{\mathrm{ico}} & 0 & 1  & 5 & 0 & 5 & 1 & 0 & 0  & 0 \\
B_3^{\mathrm{dod}} & 0 & 0  & 2 & 2 & 4 & 2 & 2 & 0  & 0 \\
B_3^{\mathrm{ico}} & 0 & 0  & 0 & 1 & 5 & 0 & 5 & 1  & 0 \\
B_4^{\mathrm{dod}} & 0 & 0  & 0 & 0 & 3 & 3 & 3 & 3  & 0 \\
B_4^{\mathrm{ico}} & 0 & 0  & 0 & 0 & 0 & 1 & 5 & 5  & 1 \\
B_5                & 0 & 0  & 0 & 0 & 0 & 0 & 0 & 12 & 0
\end{array}
\end{equation}
with integer entries, every row summing to $12$ (the degree of the
vertex graph). Collapsing to the coarser $(a, b, c)$ totals
$(|N(v) \cap \text{same cell}|, |N(v) \cap \text{next shell}|,
|N(v) \cap \text{previous shell}|)$ gives
Table~\ref{tab:P2-populations}.
\end{proposition}

\begin{table}[ht]
\centering
\caption{Coarse $(a, b, c)$ totals of the nine-cell equitable
partition of the $600$-cell vertex graph, obtained from
Equation~\eqref{eq:quotient-matrix} by collapsing within-BFS-shell
entries into $a$, outward (next BFS shell) entries into $b$, and
inward (previous BFS shell) entries into $c$. Displayed for
intuition; the proof input is the full quotient matrix
$M$.}
\label{tab:P2-populations}
\begin{tabular}{@{}lrrrrr@{}}
\toprule
Cell & Size & $a$ & $b$ & $c$ & $c/b$ \\
\midrule
$B_0$                  & $1$  & $0$  & $12$ & $0$  & --- \\
$B_1$                  & $12$ & $5$  & $6$  & $1$  & $1/6$ \\
$B_2^{\mathrm{ico}}$   & $12$ & $5$  & $6$  & $1$  & $1/6$ \\
$B_2^{\mathrm{dod}}$   & $20$ & $6$  & $3$  & $3$  & $1$ \\
$B_3^{\mathrm{ico}}$   & $12$ & $5$  & $6$  & $1$  & $1/6$ \\
$B_3^{\mathrm{dod}}$   & $30$ & $6$  & $2$  & $4$  & $2$ \\
$B_4^{\mathrm{ico}}$   & $12$ & $10$ & $1$  & $1$  & $1$ \\
$B_4^{\mathrm{dod}}$   & $20$ & $6$  & $0$  & $6$  & $\infty$ (trapped) \\
$B_5$                  & $1$  & $0$  & $0$  & $12$ & --- \\
\bottomrule
\end{tabular}
\end{table}

\begin{proof}
By exhaustive finite enumeration. For each of the $120$ vertices
$v \in V_{600}$ and each of its $12$ nearest neighbours $w \in
N(v)$, compute the BFS-shell-and-population cell of $v$ and of
$w$, and tabulate the integer count $|N(v) \cap T|$ for each
target cell $T$. This enumeration, carried out in exact integer
arithmetic in
\verb|papers/paper-v/scripts/run_rd4_exact_rational.py| (and
cross-checked in
\verb|papers/paper-v/scripts/run_rigorous_verification.py|),
yields precisely the matrix $M$ above: the same integer
$M_{S, T}$ for every $v$ in the same cell $S$. Since the count is
constant over the source cell, the partition is equitable by
\cite[\S 9.3]{GodsilRoyle2001}. The $(a, b, c)$ totals in
Table~\ref{tab:P2-populations} are obtained from $M$ by summing
entries within the source cell (giving $a$), within the next
outward BFS shell (giving $b$), and within the previous inward
BFS shell (giving $c$); all other $M_{S, T}$ entries vanish by
direct inspection.
\end{proof}

\subsection{Absorbing Markov chain on the BFS partition}
\label{sec:markov-600}

The equitable-partition property of
Proposition~\ref{prop:dual-populations} makes the nine-cell BFS
refinement of $V_{600}$ \emph{lumpable} for the uniform nearest-neighbour
random walk on the $600$-cell vertex graph in the sense of
\cite[Ch.~V]{KemenySnell1976}. We collect here the lumped
$9$-state Markov chain and the first-passage result needed by Paper
V.

\begin{definition}[Uniform nearest-neighbour walk on $V_{600}$]
\label{def:markov-600}
Let $(X_t)_{t \geq 0}$ be the discrete-time Markov chain on
$V_{600}$ that, from any vertex $v$, moves to each of its $12$
nearest neighbours with probability $1/12$. By the $2I$-vertex
transitivity (Theorem~\ref{thm:icosian-closure}) the chain is
vertex-homogeneous.
\end{definition}

\begin{proposition}[Lumped $9$-state chain, equitable lumpability]
\label{prop:lumped-markov}
Let $M \in \Z^{9 \times 9}_{\geq 0}$ be the quotient matrix of the
equitable partition of Proposition~\ref{prop:dual-populations},
displayed explicitly in Equation~\eqref{eq:quotient-matrix}. The
walk of Definition~\ref{def:markov-600} is strongly lumpable (in
the sense of \cite[Ch.~V]{KemenySnell1976}) with respect to the
nine-cell partition; its lumped process is a time-homogeneous
$9$-state Markov chain with transition matrix $\widehat P = M /
12$ (rational entries in $\{0, 1/12, 2/12, \ldots, 12/12\}$, every
row summing to $1$).
\end{proposition}

\begin{proof}
Proposition~\ref{prop:dual-populations} establishes that for every
ordered pair of cells $(S, T)$ and every vertex $v \in S$, the
number of nearest-neighbours of $v$ lying in $T$ equals the
constant $M_{S, T}$. Under the uniform nearest-neighbour walk of
Definition~\ref{def:markov-600}, the one-step transition
probability from any $v \in S$ into cell $T$ therefore equals
$M_{S, T} / 12$, independent of the choice of $v \in S$. This
source-cell-constant property is exactly the strong-lumpability
criterion of \cite[Thm.~V.5, Ch.~V]{KemenySnell1976} applied to
the uniform walk on the equitable partition
\cite[\S 9.3]{GodsilRoyle2001}. Hence the image process on the
nine-cell partition is a time-homogeneous Markov chain with
transition matrix $\widehat P = M / 12$. Row sums of $M$ are
$12$ (every vertex has degree $12$), so row sums of $\widehat P$
are $1$.
\end{proof}

\begin{theorem}[First-passage theorem $R_D(4) = 15$, exact
rational]
\label{thm:RD4-P2}
Declare the three cells $\{B_3^{\mathrm{ico}},
B_3^{\mathrm{dod}}, B_5\}$ absorbing in the $9$-state chain of
Proposition~\ref{prop:lumped-markov}, and denote by $Q \in \Q^{6
\times 6}$ the transient-to-transient block and by $R \in \Q^{6
\times 3}$ the transient-to-absorbing block of the transition
matrix. The chain is absorbing; the fundamental matrix $N :=
(I - Q)^{-1}$ and the first-passage probability matrix $B :=
N R \in \Q^{6 \times 3}$ are finite. Evaluated over $\Q$ with
Python's \texttt{fractions} module at the row corresponding to
$B_4^{\mathrm{dod}}$, one obtains the exact rational values
\[
  \Pr(\text{first hit}~B_3^{\mathrm{ico}}) = \tfrac{1}{2},
  \quad
  \Pr(\text{first hit}~B_3^{\mathrm{dod}}) = \tfrac{7}{16},
  \quad
  \Pr(\text{first hit}~B_5)                = \tfrac{1}{16},
\]
summing exactly to $1$. Therefore
\[
  \Pr(\text{first hit shell } 3) = \tfrac{15}{16},
  \qquad
  \Pr(\text{first hit shell } 5) = \tfrac{1}{16},
  \qquad
  \boxed{\;R_D(4) \;:=\;
  \Pr(\text{first hit shell } 3) / \Pr(\text{first hit shell } 5)
  \;=\; 15.\;}
\]
By the equitable-partition property, the same exact rational
values hold at each of the $20$ vertices in cell
$B_4^{\mathrm{dod}}$.
\end{theorem}

\begin{proof}
Construct $Q$ and $R$ from
Proposition~\ref{prop:lumped-markov} by restricting to the six
transient cells $\{B_0, B_1, B_2^{\mathrm{dod}},
B_2^{\mathrm{ico}}, B_4^{\mathrm{dod}}, B_4^{\mathrm{ico}}\}$;
every transient cell has a path to an absorbing cell with nonzero
probability, so the chain is absorbing and $I - Q$ is invertible
over $\Q$.

\emph{Direct derivation for the $B_4^{\mathrm{dod}}$ row.} Let
$p_{\text{ico}}$, $p_{\text{dod}}$, $p_5$ denote the first-passage
probabilities of absorption at $B_3^{\mathrm{ico}}$,
$B_3^{\mathrm{dod}}$, and $B_5$ respectively, starting from
$B_4^{\mathrm{dod}}$. From the $B_4^{\mathrm{dod}}$ and
$B_4^{\mathrm{ico}}$ rows of $\widehat P = M / 12$ in
Equation~\eqref{eq:quotient-matrix},
\[
\widehat P(B_4^{\mathrm{dod}} \to \cdot) =
\left(0, 0, 0, 0, \tfrac{3}{12}, \tfrac{3}{12},
\tfrac{3}{12}, \tfrac{3}{12}, 0\right),
\qquad
\widehat P(B_4^{\mathrm{ico}} \to \cdot) =
\left(0, 0, 0, 0, 0, \tfrac{1}{12}, \tfrac{5}{12},
\tfrac{5}{12}, \tfrac{1}{12}\right).
\]
Let $q_{\text{ico}}$, $q_{\text{dod}}$, $q_5$ denote the
analogous first-passage probabilities from $B_4^{\mathrm{ico}}$.
First-step analysis gives two coupled linear systems. From
$B_4^{\mathrm{dod}}$:
\begin{align*}
p_{\text{dod}}  &= \tfrac{3}{12} + \tfrac{3}{12}\,p_{\text{dod}}
                 + \tfrac{3}{12}\,q_{\text{dod}}, \\
p_{\text{ico}}  &= \tfrac{3}{12} + \tfrac{3}{12}\,p_{\text{ico}}
                 + \tfrac{3}{12}\,q_{\text{ico}}, \\
p_5             &= \tfrac{3}{12}\,p_5 + \tfrac{3}{12}\,q_5,
\end{align*}
and from $B_4^{\mathrm{ico}}$:
\begin{align*}
q_{\text{dod}}  &= \tfrac{5}{12}\,p_{\text{dod}}
                 + \tfrac{5}{12}\,q_{\text{dod}}, \\
q_{\text{ico}}  &= \tfrac{1}{12} + \tfrac{5}{12}\,p_{\text{ico}}
                 + \tfrac{5}{12}\,q_{\text{ico}}, \\
q_5             &= \tfrac{1}{12} + \tfrac{5}{12}\,p_5
                 + \tfrac{5}{12}\,q_5.
\end{align*}
(No terms $p_{\cdot}$-to-transient other than $B_4^{\text{dod/ico}}$
appear because every $\widehat P$-arc from $B_4^{\mathrm{dod}} \cup
B_4^{\mathrm{ico}}$ leads either back into $\{B_4^{\mathrm{dod}},
B_4^{\mathrm{ico}}\}$ or into an absorbing cell.) Eliminating the
$q$-variables: from the second pair,
$q_{\text{dod}} = \tfrac{5}{7}\,p_{\text{dod}}$,
$q_{\text{ico}} = \tfrac{1}{7} + \tfrac{5}{7}\,p_{\text{ico}}$,
$q_5 = \tfrac{1}{7} + \tfrac{5}{7}\,p_5$. Substituting into the
first,
\begin{align*}
p_{\text{dod}}  &= \tfrac{1}{4} + \tfrac{1}{4}\,p_{\text{dod}}
                 + \tfrac{1}{4}\cdot\tfrac{5}{7}\,p_{\text{dod}}
                 \;\Rightarrow\; p_{\text{dod}} = \tfrac{7}{16}, \\
p_{\text{ico}}  &= \tfrac{1}{4} + \tfrac{1}{4}\,p_{\text{ico}}
                 + \tfrac{1}{4}\bigl(\tfrac{1}{7}
                 + \tfrac{5}{7}\,p_{\text{ico}}\bigr)
                 \;\Rightarrow\; p_{\text{ico}} = \tfrac{1}{2}, \\
p_5             &= \tfrac{1}{4}\,p_5
                 + \tfrac{1}{4}\bigl(\tfrac{1}{7}
                 + \tfrac{5}{7}\,p_5\bigr)
                 \;\Rightarrow\; p_5 = \tfrac{1}{16}.
\end{align*}
These three sum to $\tfrac{7}{16} + \tfrac{1}{2} + \tfrac{1}{16}
= 1$, confirming that $\{B_3^{\mathrm{ico}}, B_3^{\mathrm{dod}},
B_5\}$ absorb all probability mass starting from
$B_4^{\mathrm{dod}}$ (so no mass leaks into $B_0$, $B_1$,
$B_2^{\mathrm{dod}}$, $B_2^{\mathrm{ico}}$). Therefore
$\Pr(\text{first hit shell }3) = p_{\text{ico}} + p_{\text{dod}}
= \tfrac{15}{16}$ and
$\Pr(\text{first hit shell }5) = p_5 = \tfrac{1}{16}$, giving
$R_D(4) = 15$ as stated.

\emph{Reproducibility check.} The same result is obtained by
computing $N = (I - Q)^{-1}$ and $B = N R$ over $\Q$ on the full
six-state transient block using Python's \texttt{fractions}
module. The script
\verb|papers/paper-v/scripts/run_rd4_exact_rational.py| carries
out this computation end-to-end and reports row values matching
the displayed rationals bit-for-bit; it is a reproducibility
artefact, not the load-bearing step of the proof.
\end{proof}

\begin{remark}[Why this is the right substrate home]
The first-passage ratio $R_D(4) = 15$ is load-bearing for the
Paper V proton-exponent corollary: replacing $R_D(4) = 15$ by the
legacy capped value $R_D(4) = 10$ shifts the proton mass ratio
from $\approx 1836$ (matching observation) to $\approx 1951$
(far off). Paper V adopts $R_D(4) = 15$ as an exact rational input
under its mass-formula ansatz; the present theorem supplies that
input as a finite-graph theorem depending only on the
equitable-partition data of
Proposition~\ref{prop:dual-populations} and the classical
lumpability / absorbing-chain theory of
\cite{KemenySnell1976, GodsilRoyle2001}, with no Paper-V-specific
hypothesis. The identification of the first-passage ratio with a
coefficient in the Paper V mass-formula exponent is a separate
model-level choice made in Paper V, not in P2.
\end{remark}

\subsection{Narrow uniqueness among the eleven regular $3$- and
$4$-polytopes}
\label{sec:polytope-uniqueness}

\begin{theorem}[Narrow $600$-cell uniqueness, computer-assisted
and conditional on imported spectral tables]
\label{thm:polytope-uniqueness}
Among the eleven regular polytopes in dimensions three and four
--- the five Platonic solids (tetrahedron, cube, octahedron,
dodecahedron, icosahedron) and the six regular $4$-polytopes
($5$-cell, $8$-cell, $16$-cell, $24$-cell, $120$-cell, $600$-cell)
--- the $600$-cell is the unique member whose vertex-graph
Laplacian spectrum contains both
\begin{enumerate}
\item[\textup{(i)}] a value in $\Q(\sqrt 5) \setminus \Q$
  (equivalently, a value involving $\varphi$), and
\item[\textup{(ii)}] the rational eigenvalue $15$.
\end{enumerate}
\end{theorem}

\begin{proof}
We check each polytope. The universal bound
$\lambda_{\max}(L) \leq 2d$ for a $d$-regular graph
\cite[Ch.~13]{GodsilRoyle2001} restricts which polytopes can carry
$\lambda = 15$. The vertex graphs of the seven polytopes with
integer Laplacian spectra --- tetrahedron ($K_4$), cube ($Q_3$),
octahedron ($K_{2,2,2}$), $5$-cell ($K_5$), $8$-cell ($Q_4$),
$16$-cell ($K_{2,2,2,2}$), and $24$-cell (spectrum
$\{0, 6, 8, 12\}$ with multiplicities $\{1, 8, 9, 6\}$) --- all
fail property~(i); their spectra are classical
\cite[Ch.~3]{BrouwerHaemers2012} and independently recomputed by
\verb|papers/paper-v/scripts/run_polytope_uniqueness_narrow.py|.
The dodecahedron (spectrum $\{0, 3 - \sqrt 5, 2, 3, 5, 3 + \sqrt
5\}$, degree $3$) and icosahedron (spectrum $\{0, 5 - \sqrt 5, 6,
5 + \sqrt 5\}$ with multiplicities $\{1, 3, 5, 3\}$, degree $5$)
both satisfy property~(i) but fail property~(ii) because
$\lambda_{\max} \leq 2d < 15$. The $120$-cell has a $4$-regular
vertex graph, so $\lambda_{\max} \leq 8 < 15$ and property~(ii)
fails. The $600$-cell satisfies both properties by
Fact~\ref{prop:spectrum-P2}. The script listed above
builds the ten lower-dimensional vertex graphs explicitly,
computes their Laplacian spectra numerically, and
verifies each of the case-checks.
\end{proof}

\begin{remark}[Scope of the uniqueness statement]
Theorem~\ref{thm:polytope-uniqueness} is confined to the eleven
regular polytopes in dimensions $3$ and $4$. In dimensions $\geq
5$ the only regular polytopes are the simplices $K_{n+1}$,
hypercubes $Q_n$, and cross-polytopes $K_{2,\ldots,2}$
\cite[Ch.~7]{CoxeterRegularPolytopes}, all of which have integer
Laplacian spectra; property~(i) therefore fails in every
dimension $\geq 5$ as well, and the statement can be extended
without providing a counterexample, though we do not promote that
extension to theorem status here. For regular polygons in
dimension $2$, the cycle graph $C_n$ has $\varphi$-valued entries
whenever $5 \mid n$ but $\lambda_{\max}(C_n) \leq 4$, so
property~(ii) always fails.
\end{remark}

\subsection{Hopf decomposition into twelve great-decagon fibres}
\label{sec:hopf-P2}

\begin{fact}[Hopf decomposition of the $600$-cell, classical
import]
\label{fact:hopf}
The $600$-cell admits a decomposition of its $120$ vertices into
$12$ disjoint great decagons (regular $C_{10}$ cycles on great
circles of $S^3$) with the nearest-neighbour adjacency of the
$600$-cell restricting to the $C_{10}$ adjacency on each fibre. The
decomposition arises from the Hopf fibration $S^3 \to S^2$
restricted to the $600$-cell
\cite[\S 22.7]{CoxeterRegularPolytopes}.
\end{fact}

\begin{proposition}[Hopf-fibre cycle spectrum]
\label{prop:hopf-spectrum}
The Laplacian eigenvalues of each $C_{10}$ fibre are
$E_k = 2 - 2\cos(2\pi k / 10)$ for $k = 0, \ldots, 9$. Explicitly:
\[
  E_0 = 0, \quad
  E_1 = 1/\varphi^2, \quad
  E_2 = 3 - \varphi, \quad
  E_3 = \varphi^2, \quad
  E_4 = 2 + \varphi, \quad
  E_5 = 4,
\]
and $E_{10-k} = E_k$ for $k = 1, \ldots, 4$. All values lie in
$\Q(\varphi)$.
\end{proposition}

\begin{proof}
Direct substitution in the closed-form cycle spectrum
$E_k = 2 - 2\cos(2\pi k/n)$ with $n = 10$, using $\cos(\pi/5) =
\varphi/2$ and $\cos(2\pi/5) = (\varphi - 1)/2 = 1/(2\varphi)$.
\end{proof}

Fact~\ref{fact:hopf} is an import from the classical $600$-cell
literature; the fibre cycle spectrum of
Proposition~\ref{prop:hopf-spectrum} is elementary given the Hopf
decomposition, and is used here only as a repository of algebraic
values for downstream citation.

%=====================================================================
\section{The $16$-Sector: $\Cl(1,3)$ as the Canonical Model}
\label{sec:16-sector}
%=====================================================================

This section presents the $16$-dimensional sector algebra. The
theorem-grade content is algebraic: $\Cl(1,3)$ is a $16$-dimensional
real associative algebra with a standard generators-and-relations
presentation, and its structure theory is classical. The
$16$-vertex tesseract plays only the role of an index set for the
monomial basis.

\subsection{Definition and basic structure}

\begin{definition}[The real Clifford algebra $\Cl(1,3)$]
\label{def:Cl}
$\Cl(1,3)$ is the real associative $\R$-algebra with unit,
generated by $\gamma_0, \gamma_1, \gamma_2, \gamma_3$ subject to
the relations
\[
  \gamma_0^2 = +1, \quad
  \gamma_1^2 = \gamma_2^2 = \gamma_3^2 = -1, \quad
  \gamma_\mu \gamma_\nu + \gamma_\nu \gamma_\mu = 0
  \text{ for } \mu \neq \nu.
\]
Equivalently, $\Cl(1,3)$ is the Clifford algebra of the quadratic
form $\operatorname{diag}(+1, -1, -1, -1)$ on $\R^4$
\cite[Ch.~16]{Lounesto}.
\end{definition}

\begin{definition}[Monomial basis]
\label{def:Cl-basis}
For each subset $A \subseteq \{0, 1, 2, 3\}$, define
\[
  \gamma_A := \gamma_{a_1} \gamma_{a_2} \cdots \gamma_{a_k},
  \qquad A = \{a_1 < a_2 < \cdots < a_k\},
\]
with $\gamma_\emptyset := 1$. The set $\{\gamma_A : A \subseteq
\{0, 1, 2, 3\}\}$ is an $\R$-basis of $\Cl(1,3)$, indexed by the
$2^4 = 16$-element power set of $\{0, 1, 2, 3\}$, which we
identify with $\Z_2^4$ via the characteristic function $A
\leftrightarrow \chi_A \in \Z_2^4$.
\end{definition}

\begin{fact}[Classical; \cite{Lounesto, Chevalley}]
\label{fact:Cl-structure}
$\dim_\R \Cl(1,3) = 16$, and $\Cl(1,3)$ is isomorphic as a real
algebra to $M_2(\HH)$, the algebra of $2 \times 2$ matrices over
the Hamilton quaternions.
\end{fact}

\begin{remark}[Tesseract as index set, not polytope structure]
\label{rmk:tesseract-index}
The index set $\Z_2^4$ above is combinatorially the vertex set of
the $4$-dimensional hypercube (tesseract). This is bookkeeping
only: the multiplication of $\Cl(1,3)$ comes from the relations
in Definition~\ref{def:Cl}, \emph{not} from any combinatorial
structure of the tesseract polytope. In particular, P2 does not
claim that the tesseract itself is an algebra or that its polytope
structure encodes Clifford multiplication.
\end{remark}

%=====================================================================
\section{The $8$-Sector: Octonions, $\mathfrak{g}_2$, and $\SU(3)$}
\label{sec:8-sector}
%=====================================================================

This section presents the octonionic $8$-sector. All theorem-grade
content is classical and cited. The algebraic core is: the octonion
algebra $\OO$; its derivation algebra $\Der(\OO) = \mathfrak{g}_2$;
and the stabilizer $\SU(3) \subset G_2 = \Aut(\OO)$ of a fixed
imaginary unit.

\subsection{The octonion algebra}

\begin{definition}[Octonions]
\label{def:octonions}
The \emph{octonion algebra} $\OO$ is the unique (up to
isomorphism) $8$-dimensional real normed division algebra. It has
a $\R$-basis $\{1, e_1, \ldots, e_7\}$ with $e_i^2 = -1$ for each
$i$, $e_i e_j = -e_j e_i$ for $i \neq j$, and multiplication
determined by the standard Fano-plane triples
\cite[\S 2.1]{BaezOctonions}.

Conjugation $x \mapsto \bar x$ is defined by $\overline{e_i} =
-e_i$ and extended $\R$-linearly; the norm is $N(x) = x \bar x
\in \R_{\geq 0}$. The imaginary part is
\[
  \Im(\OO) := \{x \in \OO : \bar x = -x\} = \R e_1 \oplus \cdots
  \oplus \R e_7, \quad \dim_\R \Im(\OO) = 7.
\]
\end{definition}

\begin{fact}[\cite{BaezOctonions, SchaferNonassociative}]
\label{fact:octonions-alternative}
$\OO$ is an alternative algebra: $(xy)x = x(yx)$ and
$x(xy) = (xx)y$ for all $x, y \in \OO$. It is non-associative:
the associator $[x, y, z] := (xy)z - x(yz)$ is not identically
zero.
\end{fact}

\subsection{The derivation algebra $\Der(\OO) = \mathfrak{g}_2$}

\begin{definition}[Derivations]
\label{def:derivations}
A \emph{derivation} of $\OO$ is an $\R$-linear map
$D \colon \OO \to \OO$ satisfying the Leibniz rule
\[
  D(xy) = D(x) y + x D(y), \qquad x, y \in \OO.
\]
The space of derivations is a real Lie algebra under the
commutator bracket $[D_1, D_2] := D_1 D_2 - D_2 D_1$; we denote
it $\Der(\OO)$.
\end{definition}

\begin{definition}[Inner-derivation map]
\label{def:Dab}
For $a, b \in \OO$, the \emph{inner derivation}
$D_{a, b} \colon \OO \to \OO$ is
\[
  D_{a, b}(x) := \big[ [a, b], x \big] - 3 [a, b, x],
\]
where $[a, b] := ab - ba$ is the commutator and $[a, b, x] :=
(ab) x - a(bx)$ is the associator.
\end{definition}

\begin{fact}[\cite{BaezOctonions, SchaferNonassociative}]
\label{fact:Der-g2}
$D_{a, b} \in \Der(\OO)$ for every $a, b \in \OO$. The derivations
$\{D_{a, b} : a, b \in \OO\}$ span $\Der(\OO)$ as a real vector
space. $\Der(\OO)$ is the compact real form of the exceptional
Lie algebra $\mathfrak{g}_2$, of real dimension $14$.
\end{fact}

\begin{remark}[Why not $\mathrm{ad}$]
\label{rmk:not-ad}
The naive map $a \mapsto \mathrm{ad}_a \colon \OO \to \OO$, $x
\mapsto [a, x]$, is not a derivation of $\OO$ in general: the
failure of the Jacobi identity in a non-associative algebra means
$\mathrm{ad}_a$ is a derivation only when $a$ is associative with
its own image. So there is no ``inclusion''
$\Im(\OO) \hookrightarrow \mathfrak{g}_2$ via $\mathrm{ad}$. The
correct map from pairs of octonions to derivations is $D_{a, b}$
of Definition~\ref{def:Dab}, which \emph{is} a derivation, by
Fact~\ref{fact:Der-g2}.
\end{remark}

\subsection{The automorphism group $G_2$ and the $\SU(3)$ stabilizer}

\begin{definition}[$G_2$]
\label{def:G2}
$G_2 := \Aut(\OO)$, the group of real algebra automorphisms of
$\OO$ (equivalently, the group of $\R$-linear maps $g \colon \OO
\to \OO$ with $g(1) = 1$ and $g(xy) = g(x) g(y)$). $G_2$ is a
compact connected real Lie group of dimension $14$, with Lie
algebra $\mathfrak{g}_2 = \Der(\OO)$ by
Fact~\ref{fact:Der-g2}.
\end{definition}

\begin{fact}[\cite{BaezOctonions}]
\label{fact:G2-14}
$G_2$ is a compact simple Lie group of dimension $14$ and rank $2$.
\end{fact}

\begin{definition}[Stabilizer of a unit imaginary octonion]
\label{def:stab-u}
Let $u \in \Im(\OO)$ with $N(u) = 1$ (so $u^2 = -1$). The
\emph{stabilizer of $u$} in $G_2$ is
\[
  \Stab_{G_2}(u) := \{g \in G_2 : g(u) = u\}.
\]
\end{definition}

\begin{fact}[$\SU(3)$ as $G_2$-stabilizer; \cite{BaezOctonions}, \cite{Macfarlane}]
\label{fact:SU3-stab}
For any unit $u \in \Im(\OO)$, the stabilizer
$\Stab_{G_2}(u)$ is isomorphic to $\SU(3)$ as a compact Lie group.
\end{fact}

\begin{remark}[Sketch of the $\SU(3)$-identification]
\label{rmk:SU3-sketch}
A fixed unit $u \in \Im(\OO)$ induces a complex structure on
$u^\perp \cap \Im(\OO) \cong \R^6$ via left multiplication by $u$:
$J_u(v) := uv$. This turns $u^\perp \cap \Im(\OO)$ into a complex
$3$-dimensional vector space. The group of automorphisms of $\OO$
fixing $u$ preserves this complex structure and also preserves the
induced Hermitian form (from the $\OO$-norm). Hence
$\Stab_{G_2}(u) \subset \mathrm{U}(3)$. Real dimension count:
$\dim G_2 - \dim \operatorname{orbit}(u) = 14 - 6 = 8 = \dim \SU(3)$,
and connectedness identifies $\Stab_{G_2}(u)$ with $\SU(3)$
\cite[\S 7]{BaezOctonions}.
\end{remark}

%=====================================================================
\section{Coefficientwise $\sigma$-Action and Compatibility}
\label{sec:sigma-compat}
%=====================================================================

This section imports the coefficientwise $\sigma$-action of
\cite{SigmaRationality} and shows that every sector-level map of
this paper defined over $\Q$ commutes with $\sigma$ after scalar
extension to $\Q(\varphi)$.

\subsection{Setting}

\begin{notation}[$\sigma$ in this paper]
\label{not:sigma-P2}
Throughout this section, $\sigma \colon \Q(\varphi) \to \Q(\varphi)$
denotes the non-trivial Galois automorphism,
$\sigma(\varphi) = 1 - \varphi$, as in
\cite[Def.~2.2]{SigmaRationality}. For any $\Q$-vector space $V$,
$\sigma_V := \id_V \otimes \sigma$ is the coefficientwise
$\sigma$-action on $V_K := V \otimes_\Q \Q(\varphi)$, per
\cite[Def.~2.5]{SigmaRationality}. The projector
$\Pi_\sigma := (\id_{V_K} + \sigma_V)/2$ is the fixed-part
projector \cite[Def.~4.1]{SigmaRationality}.
\end{notation}

\subsection{Rational forms of the sector objects}
\label{subsec:Q-forms}

To apply \cite[Thm.~4.5]{SigmaRationality} (which requires
$\Q$-linear maps between finite-dimensional $\Q$-vector spaces) we
first record explicit $\Q$-forms of each sector object.

\begin{definition}[$\Q$-forms]
\label{def:Q-forms}
\hspace{0pt}
\begin{enumerate}
\item $U_\Q := \Q^4 \subset \Q^8$, the $\Q$-form of the coordinate
  $4$-plane $U$ (Def.~\ref{def:U}).
\item $\Lambda(E_8)_\Q := \Lambda(E_8) \otimes_\Z \Q$, an
  $8$-dimensional $\Q$-vector space with basis $\Lambda(E_8)
  \otimes 1$.
\item $I_\Q := I \otimes_\Z \Q$, an $8$-dimensional $\Q$-vector
  space.
\item $\HH_\Q := \Q\{1, i, j, k\} \subset \HH$, the $\Q$-span of
  the standard quaternionic basis; $\dim_\Q \HH_\Q = 4$.
\item $\Cl_\Q(1, 3) := \Q\{\gamma_A : A \subseteq \{0,1,2,3\}\}
  \subset \Cl(1,3)$, the $\Q$-span of the monomial basis of
  Definition~\ref{def:Cl-basis}; $\dim_\Q \Cl_\Q(1,3) = 16$.
\item $\OO_\Q := \Q\{1, e_1, \ldots, e_7\} \subset \OO$, the
  $\Q$-span of the standard octonionic basis; $\dim_\Q \OO_\Q = 8$.
\end{enumerate}
For each $V \in \{U_\Q, \Lambda(E_8)_\Q, I_\Q, \HH_\Q, \Cl_\Q(1,3),
\OO_\Q\}$, write $V_K := V \otimes_\Q \Q(\varphi)$ for the scalar
extension to $\Q(\varphi)$, and $\sigma_V := \id_V \otimes \sigma$
for the coefficientwise $\sigma$-action on $V_K$
(\cite[Def.~2.5]{SigmaRationality}).
\end{definition}

\subsection{Sector-level $\sigma$-compatibility theorem}

\begin{theorem}[$\sigma$-compatibility of $\Q$-defined P2 maps]
\label{thm:sigma-compat-P2}
With the $\Q$-forms of Definition~\ref{def:Q-forms}, the following
$\Q$-linear maps each commute with $\sigma$ after scalar extension
to $\Q(\varphi)$, in the sense of
\cite[Thm.~4.5]{SigmaRationality}:
\begin{enumerate}
\item The $\Q$-linear inclusion
  $\operatorname{span}_\Q \Phi(D_4) \hookrightarrow
  \operatorname{span}_\Q \Phi(E_8) \subseteq \Lambda(E_8)_\Q$
  (well-defined since $\operatorname{span}_\Q \Phi(D_4) = U_\Q$
  and $\operatorname{span}_\Q \Phi(E_8) = \Lambda(E_8)_\Q$ by
  Proposition~\ref{prop:D4-inclusion}).
\item The $\Q$-linear isomorphism
  $(\iota_I)_\Q \colon \Lambda(E_8)_\Q \to I_\Q$ obtained by
  scalar extension of $\iota_I$
  (Fact~\ref{fact:icosian-E8}). The scalar extension to
  $\Q(\varphi)$ intertwines the icosian star-map
  (Definition~\ref{def:icosian-star}) extended to $I_K$ with the
  coefficientwise $\sigma$-action $\sigma_{\Lambda(E_8)_\Q}$ on
  $\Lambda(E_8)_K$ (Remark~\ref{rmk:star-is-sigma}).
\item The identity map $\Cl_\Q(1,3) \to \Cl_\Q(1,3)$, viewed as a
  $\Q$-linear map; its scalar extension to $\Q(\varphi)$ trivially
  commutes with $\sigma_{\Cl_\Q(1,3)}$, since the latter acts only
  on the $\Q(\varphi)$-coefficient factor.
\item The $\Q$-bilinear inner-derivation map
  $D \colon \OO_\Q \otimes_\Q \OO_\Q \to \operatorname{End}_\Q(\OO_\Q)$,
  $(a \otimes b) \mapsto D_{a, b}$, where $D_{a, b}$ is the
  $\Q$-linear endomorphism of $\OO_\Q$ defined by $x \mapsto
  D_{a,b}(x) = [[a,b], x] - 3[a, b, x]$
  (Definition~\ref{def:Dab}). $D$ is $\Q$-linear, so its scalar
  extension to $\Q(\varphi)$ commutes with $\sigma$.
\end{enumerate}
\end{theorem}

\begin{proof}
Each clause is a direct instance of
\cite[Thm.~4.5]{SigmaRationality} applied to a $\Q$-linear map
between finite-dimensional $\Q$-vector spaces:
\begin{itemize}
\item For (1), the $\Q$-linear inclusion is
  $\operatorname{span}_\Q \Phi(D_4) \hookrightarrow
  \operatorname{span}_\Q \Phi(E_8)$. Both spans are
  finite-dimensional $\Q$-vector spaces, so
  \cite[Thm.~4.5]{SigmaRationality} gives commutation with
  $\sigma$ after scalar extension.
\item For (2), $(\iota_I)_\Q := \iota_I \otimes_\Z \id_\Q$ is a
  $\Q$-linear isomorphism between $8$-dimensional $\Q$-vector
  spaces; \cite[Thm.~4.5]{SigmaRationality} gives commutation. The
  intertwining of the icosian star-map with $\sigma$ on
  $\Lambda(E_8)_K$ is Remark~\ref{rmk:star-is-sigma}: under
  $(\iota_I)_\Q \otimes \id_K$, the star-map and the
  coefficientwise $\sigma$ both act as $\varphi \mapsto
  \bar\varphi$ on the $\Q(\varphi)$-coefficient factor, so they
  coincide.
\item For (3), $\id \colon \Cl_\Q(1,3) \to \Cl_\Q(1,3)$ is
  $\Q$-linear; trivial.
\item For (4), $D$ is $\Q$-linear (the linearization of a
  $\Q$-bilinear map between finite-dimensional $\Q$-vector spaces;
  the codomain $\operatorname{End}_\Q(\OO_\Q)$ has $\Q$-dimension
  $8^2 = 64$). Apply \cite[Thm.~4.5]{SigmaRationality}.
\end{itemize}
\end{proof}

\begin{proposition}[$\eta$-intertwining; not an instance of P1]
\label{prop:eta-swap}
The map $\eta \colon I \to \HH \oplus \HH$ of
Definition~\ref{def:eta}, restricted to the $\Q$-form
$\eta_\Q \colon I_\Q \to (\HH_\Q \oplus \HH_\Q)
\otimes_\Q \Q(\varphi)$ where the codomain carries the explicit
$\Q(\varphi)$-coefficient structure of
Definition~\ref{def:real-embeddings}, satisfies
\[
  \operatorname{swap} \circ \eta_\Q
  = \eta_\Q \circ \ast,
\]
where $\operatorname{swap}(p, q) := (q, p)$ on the target and
$\ast$ is the icosian star-map. Under
Remark~\ref{rmk:star-is-sigma} (identifying $\ast$ with $\sigma$
on $I_K$), this becomes
$\operatorname{swap} \circ \eta_\Q = \eta_\Q \circ \sigma_{I_\Q}$.
This is \emph{not} the statement of
\cite[Thm.~4.5]{SigmaRationality} (which would give commutation
with the coefficientwise $\sigma$ on the codomain, not with
$\operatorname{swap}$); rather it is a separate fact about
$\eta_\Q$, which we record here for downstream use.
\end{proposition}

\begin{proof}
Let $q \in I_\Q$. By Definition~\ref{def:eta},
$\eta_\Q(q) = (q_+, q_-)$ where $q_\pm$ is obtained by applying
$\tau_\pm$ coefficientwise. Since
$\tau_+(\sigma(\lambda)) = \tau_-(\lambda)$ for every $\lambda \in
\Q(\varphi)$ (immediate from $\tau_\pm(\varphi) =
\varphi, \bar\varphi$ and $\sigma(\varphi) = \bar\varphi$), we have
$(\sigma q)_\pm = q_\mp$. Therefore
\[
  \eta_\Q(\sigma q) = ((\sigma q)_+, (\sigma q)_-) = (q_-, q_+)
  = \operatorname{swap}(q_+, q_-) = \operatorname{swap}(\eta_\Q(q)),
\]
i.e.\ $\operatorname{swap} \circ \eta_\Q = \eta_\Q \circ \sigma$
on $I_\Q$. The same identity with $\ast$ in place of $\sigma$
follows from Remark~\ref{rmk:star-is-sigma}.
\end{proof}

\begin{remark}[Why $\eta_\Q$ does not give a P1 commuting square]
\label{rmk:eta-not-P1}
P1's functoriality theorem
\cite[Thm.~4.5]{SigmaRationality} says that for a $\Q$-linear
map $T \colon V \to W$, the scalar extension $T_K$ commutes with
the coefficientwise $\sigma$ on both $V_K$ and $W_K$. For
$\eta_\Q$, the codomain $\HH_\Q \oplus \HH_\Q$ has its own
coefficientwise $\sigma$-action (from
\cite[Def.~2.5]{SigmaRationality}), which is the
$\sigma$-action on each $\HH$-factor separately. The natural map
$\eta_\Q$ does \emph{not} intertwine the source $\sigma$ with this
codomain $\sigma$ but with $\operatorname{swap}$. P1 is therefore
the wrong tool for $\eta_\Q$; Proposition~\ref{prop:eta-swap}
records the correct fact directly.
\end{remark}


\begin{remark}[Why $\pi_H$ does not have a $\sigma$-equivariant
form]
\label{rmk:pi-H-no-sigma}
The composition $\pi_H = \operatorname{pr}_{\rm phys} \circ \eta
\circ \iota_I$ projects out the $\tau_-$-coefficient information
that $\sigma$ would have permuted with the $\tau_+$-coefficient.
Hence any putative codomain $\sigma$-action on $\HH$ alone is
not canonically determined by $\sigma$ on the source $I_K$, and
the diagram for $\pi_H$ \emph{by itself} does not commute with any
non-trivial codomain involution. This is correct and intentional:
$\pi_H$ is a strictly one-sided projection by design. Downstream
papers that need $\sigma$-equivariance should use
$\eta_\Q$ (which has the swap intertwining) rather than $\pi_H$
(which does not).
\end{remark}

%=====================================================================
\section{Citation Contract and Downstream Handoff}
\label{sec:handoff}
%=====================================================================

\subsection{What downstream papers may cite}

\begin{enumerate}
\item (P3, P5, Foundations.)
  Proposition~\ref{prop:D4-inclusion}: $\Phi(D_4) = \Phi(E_8)
  \cap U$, i.e.\ the strict sub-root-system inclusion $D_4 \subset
  E_8$ via the coordinate $4$-plane.
\item (P3, P5, Foundations.) Definition~\ref{def:V24} and
  Proposition~\ref{prop:24-in-600}: the explicit $24$-cell vertex
  set $V_{24} \subset V_{600}$ with its isometry to the $D_4$ root
  polytope.
\item (P3, P10, Foundations.)
  Definitions~\ref{def:icosian-ring}--\ref{def:pi-H} and
  Theorem~\ref{thm:pi-H}: the icosian ring, star-map, and the
  $E_8 \to H_4$ projection with its densification property and its
  recovery of the $H_4$ root system.
\item (Paper V, Paper XXII, P3, Foundations.)
  Theorem~\ref{thm:icosian-closure}: the $600$-cell vertex set
  $V_{600}$ is a subgroup of $\mathrm{Sp}(1)$ isomorphic to the
  binary icosahedral group $2I$ of order $120$, verified in
  exact $\Q(\sqrt 5)$-arithmetic.
\item (Paper V, Paper XXII, P3, Foundations.)
  Theorem~\ref{thm:shell-class}: the nine Euclidean distance
  shells from the identity (Table~\ref{tab:P2-shells}) are in
  size-preserving bijection with the nine $2I$-conjugacy classes,
  with each shell equal to a single class.
\item (Paper XXII, P3, Foundations.)
  Corollary~\ref{cor:class-sums-central}: under the isomorphism
  $\C^{V_{600}} \cong \C[2I]$ afforded by the left-regular
  action, each shell-adjacency operator $\mathbf{A}_k$ is
  right-convolution by the central class sum $C_k \in Z(\C[2I])$,
  and therefore acts by a scalar on every isotypic component of
  $\C[2I]$.
\item (Paper V, Paper XXII, Foundations.)
  Fact~\ref{prop:spectrum-P2} and
  Corollary~\ref{cor:94-26-split}: the Laplacian spectrum of the
  $600$-cell vertex graph with its isotypic multiplicities, and
  the resulting $W_{\mathrm{int}} \oplus W_{\mathrm{irr}} = 94 +
  26$ integer / irrational decomposition preserved by every
  polynomial in the shell-adjacency matrices. Downstream papers
  citing this item inherit the imported $2I$ character-table
  dependence recorded in \S\ref{sec:intro}; the invariance of the
  decomposition follows from
  Corollary~\ref{cor:class-sums-central} and is independent of
  the character-table step.
\item (Paper V, P3.) Proposition~\ref{prop:BFS-shells} and
  Proposition~\ref{prop:dual-populations}: the six BFS shells
  with sizes $\{1, 12, 32, 42, 32, 1\}$, the icosahedral /
  dodecahedral subpopulation split at shells $2, 3, 4$, and the
  equitable-partition property of the resulting nine-cell
  partition.
\item (Paper V, Foundations.)
  Proposition~\ref{prop:lumped-markov} and
  Theorem~\ref{thm:RD4-P2}: the $9$-state lumped Markov chain on
  the BFS partition with integer out-edge counts per $12$, and
  the exact rational first-passage identity $R_D(4) = 15$
  obtained by solving $(I - Q)^{-1} R$ over $\Q$ with shells $3$
  and $5$ absorbing.
\item (Paper V, Foundations.)
  Theorem~\ref{thm:polytope-uniqueness}: the $600$-cell is the
  unique member of the eleven-polytope family whose Laplacian
  spectrum carries both a $\Q(\sqrt 5) \setminus \Q$ value and
  the rational $15$.
\item (Paper V, Foundations.) Fact~\ref{fact:hopf} and
  Proposition~\ref{prop:hopf-spectrum}: the Hopf decomposition
  into twelve great-decagon fibres with its $C_{10}$
  cycle-spectrum values $\{0, 1/\varphi^2, 3-\varphi, \varphi^2,
  2+\varphi, 4\}$.
\item (Foundations.) Fact~\ref{fact:Cl-structure}: the algebraic
  $16$-sector $\Cl(1,3) \cong M_2(\HH)$ with its $\Z_2^4$ monomial
  basis.
\item (P8, P9, Foundations.)
  Definitions~\ref{def:octonions}--\ref{def:stab-u} and
  Facts~\ref{fact:Der-g2}, \ref{fact:SU3-stab}: the octonionic
  $8$-sector, the derivation algebra $\Der(\OO) =
  \mathfrak{g}_2$ via $D_{a, b}$, and the stabilizer $\SU(3)
  \subset G_2$.
\item (Foundations, all branches.) Theorem~\ref{thm:sigma-compat-P2}:
  every $\Q$-defined P2 map commutes with the coefficientwise
  $\sigma$-action after scalar extension to $\Q(\varphi)$, per
  \cite{SigmaRationality}.
\end{enumerate}

\subsection{Recommended citation sentence}

``By P2, the cascade algebraic substrate consists of:
an explicit strict $D_4 \subset E_8$ sub-root-system inclusion
(Proposition~\ref{prop:D4-inclusion}); an icosian-based projection
$\pi_H$ from $E_8$ lattice data to $H_4$ data
(Theorem~\ref{thm:pi-H}); an explicit $24$-cell inside the
$600$-cell (Proposition~\ref{prop:24-in-600}); the identification
of the $600$-cell vertex set with the binary icosahedral group
$2I$ under Hamilton multiplication and the shell-to-conjugacy-class
bijection (Theorems~\ref{thm:icosian-closure}--\ref{thm:shell-class})
with the associated $94{+}26$ isotypic decomposition of
$\C^{V_{600}}$ (Corollary~\ref{cor:94-26-split}); the BFS shell
structure with its equitable icosahedral / dodecahedral
subpopulations (Propositions~\ref{prop:BFS-shells}--\ref{prop:dual-populations});
the narrow uniqueness of the $600$-cell among the eleven regular
$3$-/$4$-polytopes (Theorem~\ref{thm:polytope-uniqueness}); the
Hopf decomposition into twelve great-decagon fibres with their
$C_{10}$ cycle spectrum (Proposition~\ref{prop:hopf-spectrum});
the real Clifford algebra $\Cl(1,3)$ as the $16$-sector model
(Fact~\ref{fact:Cl-structure}); the octonionic $8$-sector with
$\Der(\OO) = \mathfrak{g}_2$ (Fact~\ref{fact:Der-g2}) and
$\SU(3) = \Stab_{G_2}(u)$ (Fact~\ref{fact:SU3-stab});
and the coefficientwise $\sigma$-compatibility of all $\Q$-defined
maps (Theorem~\ref{thm:sigma-compat-P2}).''

\subsection{What P2 does not provide}

\begin{enumerate}
\item No seven-rung uniqueness theorem. The cascade chain
  $E_8 \to H_4 \to 40 \to D_4 \to 16 \to 8 \to 0$ is not proved
  to be forced by any external principle in this paper.
\item No crystallographic inclusion of $H_4$ into $E_8$. $H_4$
  appears only as the image of the projection $\pi_H$.
\item No $2I/2T$ quotient structure on $V_{600}$. $V_{24}$ is an
  inscribed subset, not a fibre quotient.
\item No refinement-graph construction. Refinement graphs are the
  subject of P3.
\item No closure-functional / dynamics content. That is the
  subject of P4.
\item No gauge, spectral, hydrodynamic, or computational
  correspondence. Those are the subjects of downstream branch
  papers.
\item No physical identification of the sectors with gauge fields,
  metrics, spinors, or any specific Millennium-problem target.
\end{enumerate}

\appendix

%=====================================================================
\section{Notation Audit}
\label{app:notation}
%=====================================================================

\begin{center}
\begin{tabular}{lll}
\hline
Symbol & Meaning & Defined \\
\hline
$\Phi(E_8)$ & $E_8$ root system in $\R^8$ & Def.~\ref{def:E8} \\
$\Lambda(E_8)$ & $E_8$ root lattice & Def.~\ref{def:lambda-E8} \\
$U$ & coordinate $4$-plane $\operatorname{span}(e_1, \ldots, e_4)$ & Def.~\ref{def:U} \\
$\Phi(D_4)$ & $D_4$ root system in $U$ & Def.~\ref{def:D4} \\
$\HH$ & Hamilton quaternions & notation \\
$I$ & icosian ring & Def.~\ref{def:icosian-ring} \\
$\ast$ & icosian star-map $\varphi \mapsto \bar\varphi$ & Def.~\ref{def:icosian-star} \\
$\iota_I$ & lattice isometry $\Lambda(E_8) \to I$ & Fact~\ref{fact:icosian-E8} \\
$\tau_\pm$ & real embeddings of $\Q(\varphi)$ & Def.~\ref{def:real-embeddings} \\
$\eta$ & twin map $I \to \HH \oplus \HH$ & Def.~\ref{def:eta} \\
$\operatorname{pr}_{\rm phys}$ & physical-space projection $\HH \oplus \HH \to \HH$ & Def.~\ref{def:pr-phys} \\
$\pi_H$ & $E_8 \to H_4$ projection & Def.~\ref{def:pi-H} \\
$V_{600}, V_{24}$ & $600$-cell, inscribed $24$-cell vertex sets & Defs.~\ref{def:V600}, \ref{def:V24} \\
$\Cl(1,3)$ & real Clifford algebra of signature $(1,3)$ & Def.~\ref{def:Cl} \\
$\gamma_\mu, \gamma_A$ & Clifford generators / monomial basis & Defs.~\ref{def:Cl}, \ref{def:Cl-basis} \\
$\OO$ & octonion algebra & Def.~\ref{def:octonions} \\
$\Im(\OO)$ & imaginary octonions & Def.~\ref{def:octonions} \\
$D_{a,b}$ & inner-derivation map on $\OO$ & Def.~\ref{def:Dab} \\
$\Der(\OO)$ & derivation algebra of $\OO$ = $\mathfrak{g}_2$ & Def.~\ref{def:derivations}, Fact~\ref{fact:Der-g2} \\
$G_2$ & automorphism group $\Aut(\OO)$ & Def.~\ref{def:G2} \\
$\Stab_{G_2}(u)$ & stabilizer of unit $u \in \Im(\OO)$, = $\SU(3)$ & Def.~\ref{def:stab-u}, Fact~\ref{fact:SU3-stab} \\
$\sigma$ & coefficientwise Galois action $\varphi \mapsto \bar\varphi$ & Notation~\ref{not:sigma-P2}, per \cite{SigmaRationality} \\
\hline
\end{tabular}
\end{center}

\begin{thebibliography}{99}

\bibitem{SigmaRationality}
\emph{Cascade--Classical Correspondence I: Sigma-Fixed Rationality
over $\Q(\varphi)$},
\texttt{papers/cascade-sigma-rationality/cascade-sigma-rationality.tex}.
(Paper P1 of the cascade infrastructure programme.)

\bibitem{CascadeInfrastructurePlan}
\emph{Cascade Infrastructure Plan},
\texttt{papers/cascade-infrastructure-plan.md}.

\bibitem{Humphreys}
J.~E.~Humphreys, \emph{Introduction to Lie Algebras and
Representation Theory}, Graduate Texts in Mathematics \textbf{9},
Springer, 1972. The standard model of $D_4$ root system used here
(\S\ref{sec:D4}) appears as part of the classical-type root system
catalogue in \S 12.1 (after the general theory developed in
Ch.~III). The crystallographic restriction in $\R^n$ used in
Lemma~\ref{lem:H4-noncrystallographic} is Theorem 6.4 (in some
printings, Theorem 6.4 of the chapter on root systems).

\bibitem{Bourbaki}
N.~Bourbaki, \emph{Groupes et alg\`ebres de Lie, Chapitres 4, 5 et
6}, Masson, 1981 (Hermann, 1968 original). Root systems of type
$E_8$: Planches.

\bibitem{FultonHarris}
W.~Fulton and J.~Harris, \emph{Representation Theory: A First
Course}, Graduate Texts in Mathematics \textbf{129}, Springer,
1991. Triality and $\Spin(8)$: \S 20.

\bibitem{Serre}
J.-P.~Serre, \emph{A Course in Arithmetic}, Graduate Texts in
Mathematics \textbf{7}, Springer, 1973. The $E_8$ lattice and its
basic structural properties (even, unimodular, self-dual) appear in
Ch.~V (Quadratic forms with discriminant $\pm 1$); the standard
explicit description used in \S\ref{sec:E8} is in \S V.1.4 (or
\S III.4 in the parallel discussion of the discriminant of the
$E_8$ form).

\bibitem{CoxeterRegularPolytopes}
H.~S.~M.~Coxeter, \emph{Regular Polytopes}, 3rd ed., Dover, 1973.
$H_4$ Coxeter group: \S 3; $600$-cell coordinates: \S 22.7;
$24$-cell: \S 3.3.

\bibitem{ConwaySloane}
J.~H.~Conway and N.~J.~A.~Sloane, \emph{Sphere Packings, Lattices
and Groups}, 3rd ed., Springer, 1999. Icosians and $E_8$ lattice:
Ch.~8.

\bibitem{MoodyPatera}
R.~V.~Moody and J.~Patera, \emph{Quasicrystals and icosians}, J.
Phys.\ A: Math.\ Gen.\ \textbf{26} (1993), 2829--2853. Explicit
icosian / $H_4$ / $E_8$ interface.

\bibitem{Lounesto}
P.~Lounesto, \emph{Clifford Algebras and Spinors}, 2nd ed., London
Math.\ Soc.\ Lecture Note Series \textbf{286}, Cambridge
University Press, 2001. Real Clifford algebras and the matrix
classification: Ch.~16 ($\Cl_{p,q}$ as a matrix algebra over
$\R, \C, \HH$); the specific identification $\Cl(1,3) \cong
M_2(\HH)$ appears in Table~16.4 (with the equivalent notation
$\mathrm{Cl}_{1,3} \cong M(2, \HH)$).

\bibitem{Chevalley}
C.~Chevalley, \emph{The Algebraic Theory of Spinors and Clifford
Algebras}, Collected Works Vol.~2, Springer, 1997. Clifford
algebras over $\R$.

\bibitem{BaezOctonions}
J.~Baez, \emph{The Octonions}, Bull.\ AMS \textbf{39} (2002),
145--205. Octonions, $\Der(\OO) = \mathfrak{g}_2$, $G_2$, and
$\SU(3) = \Stab_{G_2}(u)$: \S\S 2, 4, 7.

\bibitem{SchaferNonassociative}
R.~D.~Schafer, \emph{An Introduction to Nonassociative Algebras},
Dover, 1995 (Academic Press 1966 reprint). $\Der(\OO) =
\mathfrak{g}_2$: Ch.~III.

\bibitem{Macfarlane}
A.~J.~Macfarlane, \emph{The sphere $S^6$ viewed as a $G_2/\SU(3)$
coset space}, Int.\ J.\ Mod.\ Phys.\ A \textbf{17} (2002),
2595--2613. $\SU(3)$ stabilizer in $G_2$.

\bibitem{DuVal1964}
P.~Du~Val, \emph{Homographies, Quaternions and Rotations}, Oxford
University Press, 1964. Classical classification of finite
subgroups of $\Sp(1)$; the binary polyhedral groups (Ch.~3).

\bibitem{BrouwerHaemers2012}
A.~E.~Brouwer and W.~H.~Haemers, \emph{Spectra of Graphs},
Springer, 2012. Laplacian spectra of Platonic and regular-polytope
vertex graphs (Ch.~3); spectra of Cayley-graph and
distance-regular graph constructions.

\bibitem{GodsilRoyle2001}
C.~Godsil and G.~Royle, \emph{Algebraic Graph Theory}, Graduate
Texts in Mathematics~\textbf{207}, Springer, 2001. Equitable
partitions and quotient matrices (\S 9.3); Laplacian spectra and
bounds (Ch.~13).

\bibitem{KemenySnell1976}
J.~G.~Kemeny and J.~L.~Snell, \emph{Finite Markov Chains}, 2nd
ed., Springer, 1976. Lumpability theorem and absorbing chains
(Ch.~V); fundamental matrix $N = (I - Q)^{-1}$ for absorbing
chains.

\bibitem{CascadeFoundations}
\emph{Cascade Foundations: F1--F8} (internal note, formal
framework), \texttt{papers/cascade-derivation/cascade-foundations.md}.
Excluded from theorem support per \S\ref{sec:intro}.

\bibitem{CascadePregeometry}
\emph{Cascade Pregeometry: P0--P9} (internal note, formal
framework), \texttt{papers/cascade-derivation/cascade-pregeometry.md}.
Excluded from theorem support per \S\ref{sec:intro}.

\bibitem{CascadeEmbeddings}
\emph{Cascade Embeddings} (internal note),
\texttt{papers/cascade-derivation/cascade-embeddings.md}.
Historical motivation only. \S 12 (computational verifications)
is not used as theorem-grade support.

\end{thebibliography}

\end{document}
